\section{Heegner points and transfer}
\label{sec:heegner-interpolation}

We prove the Heegner comparison, the genus and ring bounds, and the
interpolation criterion stated in Section~\ref{sec:transfer-criteria}.
The point and its paired functional define one determinant coordinate.
Its central value measures the point index after the Selmer and local
contributions have been subtracted. Local switches then control its
divisors. An outer limit supplies the uniform ring bound when residual
irreducibility is unavailable; a universal binary family supplies
interpolation.

\subsection{Height and base Selmer comparisons}

Retain the genus and ring data, character sums $P_c(e)$, and lattice
indices of Section~\ref{sec:transfer-criteria}. We first justify the two
comparisons used there.

\begin{proof}[Proof of Proposition~\ref{prop:heegner-comparison}]
The explicit Gross--Zagier formula for ring classes
\citep[Theorem~1.1, pp.~2524--2525]{CaiShuTian2014} applies. Automorphic induction of the
quadratic character gives the product in Equation~\eqref{eq:h:height}. Discriminant and conductor are prime to $N$, and the
classical split Heegner hypothesis holds; consequently the shared-prime
corrections in the general formula are absent. The unit-index convention
is constant in the fields retained here. With a character sum, solving
the formula for the height gives exactly $c\sqrt{|k|}$, as in
Equation~\eqref{eq:h:height}. The forward Gross--Zagier--Kolyvagin theorem
gives the rank and finiteness assertion
\citep{GrossZagier1986,Kolyvagin1990}.

For Equation~\eqref{eq:h:index-comparison}, both free points are compared
with the free line on $E^{(h)}(\Q)$. Its regulator cancels from the ratio
of their heights, since absolute canonical height is unchanged by twist
identification. This height invariance follows from the canonical-height
limit and the degree formula for heights under morphisms
\citep[Chapter~IV, Lemmas~4.1 and~4.5--4.6, and Remark~4.12]{MilneEC2021}.
The remaining variable factor is
$\sqrt{|hk|}L(E^{(hk)},1)$; its infinity sign is that of $h_*$. Period
and lattice changes have bounded $2$-valuation. If the rank-one line is
instead that of the fixed curve $E^{(h_*)}$, its regulator is fixed and
the same calculation gives Equation~\eqref{eq:h:reverse-comparison}.
\end{proof}

\begin{proof}[Proof of Lemma~\ref{lem:h:base-selmer-bound}]
Restriction, twisting, and addition or subtraction of conjugates give
maps in both directions between the Selmer group over $K$ and
\[
 \Sel_{2^\infty}(E^{(h_0)}/\Q)
 \oplus\Sel_{2^\infty}(E^{(h_0k)}/\Q)
\]
whose composites are multiplication by $2$. They induce isogenies on
divisible subgroups. On the finite quotients the kernel of either induced
map is killed by $2$: subtract an element of the divisible subgroup to
lift a divisible image, and then use the composite. The generator
dimensions are uniformly bounded. Indeed the mod-$2$ Selmer group over
$K$ is bounded by Kummer cohomology over $K(E[2])$; this is a quadratic
extension of a fixed field with a bounded number of ramified primes.
Kummer theory bounds the generator number of the corresponding pro-$2$
group over the fixed field, and the index-two subgroup bound preserves
a uniform bound. The fixed local support includes that of $h_0$.
The finite quotient on the displayed product is bounded by hypothesis;
therefore its counterpart over $K$ has bounded length as well. For a
finite collection of fields, finiteness of each quotient is enough.
\end{proof}

\subsection{Integral models and ring-class deformation}
\label{rl:section}\label{sec:h:finite-models}

The determinant argument uses characters of increasing conductor.
We construct their cochain limits together with maps to the original
cochains. These maps will let us detect limiting Selmer classes by
Galois evaluation and choose local switches. Later we retain the same
maps through the outer limit needed for the ring bound.

Write $T=T_2E$ and $\Lambda=\mathbf Z_2[[t]]$. For a quadratic
squareclass $h$, put $T^h=T\otimes\chi_h$.  The imaginary quadratic field
$K$ satisfies the split Heegner hypothesis at $2N$.  We omit the two fields
with extra units and, when $E$ has complex multiplication, its CM field.
A quadratic character $\chi$ comes from a quadratic twist over $\mathbf Q$.
We use
\[
 A_i=(\mathbf Z/2^i\mathbf Z)[t]/(t^i).
\]
The same constructions work over $A_i[G]$ and $\Lambda[G]$, where $G$ is a
fixed finite elementary abelian $2$-group and the additional character has
values in $G$.  Statements uniform in the number of generators of $G$ will
be identified separately; finite-diagram compactness by itself is a statement
for fixed $G$.

\subsubsection{Finite models with their cochain maps}

The general finite-model construction is Lemma~\ref{rl:finite-model}.
We apply it to the following arithmetic cochains and retain its actual
contractions throughout.

For a profinite group and a finite free coefficient module over an Artin local ring $A$,
continuous inhomogeneous cochains are flat: locally constant functions are
the filtered union of the modules of functions on finite clopen partitions,
and these modules are finite free.  Their reduction is the cochain complex
with reduced coefficients.  The same observations apply to restrictions,
mapping cones, and local conditions given by finite free complexes.  In the
applications below the groups are Galois groups of totally imaginary number
fields or nonarchimedean local fields.  Their $2$-cohomological dimension and
the usual finiteness theorems give the bounded residual cohomology required
in Lemma~\ref{rl:finite-model}.

For an inner limit, the quadratic field $K$, the active conductor support,
the finite group $G$, and the number of derivative primes are fixed.
The residual representation is $E[2]$: all the scalar characters and
group-like coefficients become trivial in the residue field. Put
$F=K(E[2])$. The number of places of $F$ in the permitted set $S_i$ is
bounded even when its auxiliary primes move. Kummer theory bounds
$H^1(G_{F,S_i},\mathbf F_2)$ by the $S_i$-units modulo squares and
$\operatorname{Cl}(F)[2]$; their dimensions are bounded by the unit
rank, $\#S_i$, and the fixed class-group dimension. Hochschild--Serre
for the fixed group $\operatorname{Gal}(F/K)$ then bounds the residual
global $H^1$ for $E[2]$. The global Euler characteristic bounds $H^2$,
and $H^0$ is already bounded. Local cohomology and the specified local
conditions have bounded residual dimensions, since the local degrees
of $K$ are at most two and the number of places is bounded. Thus the
target Selmer cones over $K$ have uniformly bounded finite-model ranks
in this inner limit. This bound may depend on the fixed active support;
the stronger outer bound is proved in Lemma~\ref{rl:residual-bound}.
No rank bound is asserted for cochains over the growing ring-class
fields, or for the induced modules in their Shapiro descriptions.

\begin{lemma}[Limits of finite diagrams]\label{rl:finite-limits}
Suppose diagrams of bounded finite free complexes over $A_i$, or
over $A_i[G]$, have a common bounded cohomological degree interval
and bounded ranks in every degree. Include the chain homotopies
and inverses for every asserted identity or equivalence. Along a
cofinal subsequence, or a nonprincipal ultrafilter, their matrices
converge to the same diagram of perfect complexes over $\Lambda$
or $\Lambda[G]$.

Moreover, the contractions into the original cochains may be retained for
the evaluation construction in Lemma~\ref{rl:evaluation}.  Adding finitely
many new objects or maps does not require changing already chosen models
for the old objects. Finite-stage arithmetic operations with unbounded
source complexes are permitted when only finitely many of their outputs,
and the witnesses to their identities, are retained in the bounded
target complexes.
\end{lemma}

\begin{proof}
There are finitely many possible graded ranks.  Fix them on an
ultrafilter-large set.  At a fixed precision $j$, all matrix entries belong
to the finite ring $A_j$ (or $A_j[G]$), so have unique ultrafilter limits.
These limits are compatible as $j$ varies.  Matrix identities, including
the homotopy-inverse identities, are finite polynomial identities and pass
to the limit.  New models and maps may be mapped to old models by their
already fixed contractions.  Passing to any slower cofinal precision has
the same old limits.  For the all-sequence evaluation statements below,
fix this ultrafilter once and for all, and keep the original contraction
at every stage.  Later arithmetic diagrams must extend these same
stages.  Restriction to an ultrafilter-large set and passage to slower
cofinal coefficient precision preserve every old evaluation on every
old group sequence. Thus later operations use the same stage diagrams
and ultrafilter, preserving the arithmetic maps already constructed.

Here is the last assertion in cochain terms. Let $D_i$ be an original
target cochain complex, with contraction
$(\iota_i,\pi_i,h_i)$ onto its bounded finite model $C_i$.
Perform transfer and scalar weighting at stage $i$ before applying
this contraction. A resulting cycle $z_i'\in D_i^1$ gives
$z_i=\pi_i z_i'\in C_i^1$, and
\[
 z_i'-\iota_i z_i=d(h_i z_i').
\]
Thus its relation to the original arithmetic cocycle is retained by
the original contraction. If an identity in the target is
$u_i'-v_i'=d w_i'$, retain the vectors
$\pi_i u_i',\pi_i v_i',\pi_i w_i'$; the corresponding identity is
still exact. For identities involving maps between target complexes,
use the transported maps $\pi_i'f_i\iota_i$ and their contraction
homotopies, so that the same assertion includes the required boundary
corrections. Local lifts, closed degree-two functionals, and witnesses
to specialization, duality and switch identities are treated in the
same way. They are finitely many vectors or matrices of bounded size.
Their entries, as well as the differentials, therefore have simultaneous
limits. Neither a free model of the arithmetic source nor the entire
transfer map from that source is included in the limiting diagram.
Any precision lost in a finite-stage calculation is accounted for
before retaining these vectors; passing to the same slower cofinal
precision for the whole diagram preserves all its identities.
\end{proof}

The finite models now retain both the original cochains and their
comparison maps. We next specify the local conditions on those cochains.
The point is to retain the actual Kummer images at every place dividing $2N$;
the determinant comparisons will use their integral orthogonality, not
just the rational Selmer ranks.

\subsubsection{Integral local conditions and duality}

\begin{lemma}[Integral Kummer local condition]\label{rl:kummer}
\label{lem:h:duality}
Let $F$ be a finite extension of a nonarchimedean local field of
characteristic zero, and let $B/F$ be any elliptic curve.  Put
$M_F=B(F)^\wedge_2=\varprojlim_n B(F)/2^nB(F)$.  There is a perfect local
condition
\[
 U_F=M_F[-1]\longrightarrow R\Gamma(F,T_2B).
\]
For every $n\geq1$, derived reduction has
\[
 H^0(U_F\otimes^{\mathbf L}\mathbf Z/2^n)=B(F)[2^n],\qquad
 H^1(U_F\otimes^{\mathbf L}\mathbf Z/2^n)=B(F)/2^nB(F).
\]
Its map in degree zero is the isomorphism onto
$H^0(F,B[2^n])$, and its map in degree one is the finite Kummer injection.
The condition is its own exact orthogonal complement under local Tate
duality, for every reduction type and with no hypothesis on $B[2]$.
\end{lemma}

\begin{proof}
The local analytic description of $B(F)$ shows that $M_F$ is a finitely
generated $\mathbf Z_2$-module.  Its torsion is the finite group
$B(F)[2^\infty]$.  Consequently $M_F$ has a free resolution of length at
most one.  Also $H^0(F,T_2B)=0$, since $B(F)[2^\infty]$ has bounded
exponent.  The integral Kummer injection into $H^1(F,T_2B)$ therefore
defines the displayed morphism using the truncation in degrees at most
one of $R\Gamma(F,T_2B)$.

Tensor the free resolution of $M_F$ with $\mathbf Z/2^n$.  Its two
cohomology groups are $M_F[2^n]$ and $M_F/2^nM_F$. The inverse limits
of the ordinary Kummer sequences identify their maps as asserted.

For orthogonality, compose cup product and the Weil pairing with the local
invariant map.  A morphism
\[
 M_F[-1]\otimes^{\mathbf L}M_F[-1]
       \longrightarrow\mathbf Z_2[-2]
\]
is determined by its homomorphism $M_F\otimes M_F\to\mathbf Z_2$:
the possible $\operatorname{Tor}_1$ is in degree one and gives no extra
degree-zero morphism to this target.  The indicated homomorphism vanishes,
because the finite Kummer image annihilates itself at every level by local
Tate duality.  Choose a nullhomotopy in free models.  The induced morphism
\[
 R\Gamma(F,T_2B)/U_F
       \longrightarrow R\operatorname{Hom}_{\mathbf Z_2}
                         (U_F,\mathbf Z_2)[-2]
\]
is a quasi-isomorphism.  This can be checked after reduction modulo $2$:
degree one is the perfect pairing of the finite Kummer quotient with the
Kummer image, and degree two is dual to $B(F)[2]$.  Its perfect cone is
therefore zero by derived Nakayama.  This proves exact orthogonality and
also its compatibility with derived reduction.  This is precisely the
finite-module duality input of \cite[Proposition~5.2.4]{Nekovar2006}.
\end{proof}

Where the scalar deformation is trivial, extend this local condition by
derived scalar extension to $A_i$ or $A_i[G]$.  At an unramified good prime
use the inflated residue-field complex
\[
 U_v=[\,M\xrightarrow{\operatorname{Fr}_v-1}M\,]
 \quad\text{in degrees }0,1.
\]
Its quotient in full local cochains is the singular local complex.
The full local condition is the identity on local cochains, and its
orthogonal strict condition is zero.  At a derivative prime the
conditions used later are the following specified truncations instead.
In its split local model write
\[
 L_v=D^0\oplus(F_v\oplus S_v)[-1]\oplus D^2[-2],
\]
where the differential is zero, $D^0$ is dual to $D^2$, and the finite
and transverse planes $F_v,S_v$ are each isotropic.  Put
\[
 U_{\rm str}=D^0,\qquad
 U_{\rm rel}=D^0\oplus(F_v\oplus S_v)[-1],\qquad
 U_F=D^0\oplus F_v[-1],\qquad U_S=D^0\oplus S_v[-1].
\]
Then $U_{\rm str}^{\perp}=U_{\rm rel}$ and each plane condition is
its own orthogonal complement.  Every displayed condition maps
isomorphically onto local $H^0$ and injectively into local $H^1$.
Thus they satisfy precisely the detection hypotheses used below.
Their finite arithmetic realization and the isotropic-plane
nullhomotopies are part of the derivative construction; these
truncated conditions are not identified with the full local condition.

For a finite permitted set $S$ the Selmer complex is, by definition,
\begin{equation}\label{rl:selmer-definition}
 C_S=\operatorname{Cone}\left(
 R\Gamma(G_{K,S},M)\oplus\bigoplus_{v\in S}U_v
 \longrightarrow\bigoplus_{v\in S}R\Gamma(K_v,M)
 \right)[-1].
\end{equation}
At $2N$ take the Kummer conditions just constructed.  Restriction and
localization triangles are the triangles of this cone.  The finite rings
$A_i$ and $A_i[G]$ are Frobenius coefficient rings; alternatively the
duality maps may be checked after reduction to their residue field.
Exact local orthogonality and global duality therefore give
\begin{equation}\label{rl:global-duality}
 C_S\simeq R\operatorname{Hom}(C_S^{\perp},A_i)[-3],
\end{equation}
with the appropriate dual representation.  This application uses the
orthogonality hypotheses of
\cite[Theorem~6.3.4]{Nekovar2006}; the theorem is not being used to assert
the existence of a varying-prime limit.  Complex places contribute zero
Tate complexes.  If the character is inverted by the nontrivial
automorphism of $K/\mathbf Q$, conjugate transport followed by the Weil
pairing identifies the two representations over the same scalar ring.
All maps and nullhomotopies in this paragraph are finite diagrams to
which Lemma~\ref{rl:finite-limits} applies.

\subsection{The moving conductor character}
\label{sec:h:conductor-direction}

The conductor character must preserve the local conditions at $2N$
while making the new local terms invertible after reduction modulo $2$
and inversion of $t$. The following proposition supplies those local
values. Its separate bounds on Tate Frobenius control the contribution
of the moving prime when we later specialize to $t=0$.

\begin{proposition}[Ring-class characters with prescribed local values]
\label{rl:characters}\label{lem:h:conductor-primes}
Fix finitely many split primes of $K$ above odd rational primes, outside
$2N$, and orient one prime in each conjugate pair.  There are distinct
split rational primes $r_i$ avoiding all the fixed data and characters
\[
 \xi_i:\operatorname{Gal}(H_{r_i}/K)\longrightarrow
                 \mathbf Z/2^{n_i}\mathbf Z,
 \qquad n_i\longrightarrow\infty,
\]
such that $\xi_i$ is surjective on the inertia quotient at $r_i$, is zero
at the primes above $2N$, and has odd value at each prescribed oriented
prime.  We can require $r_i\to1$ in $\mathbf Z_2$ and
\begin{equation}\label{eq:h:r-avoidance}
 v_2(2+a_{r_i}(E)),\quad v_2(2-a_{r_i}(E)),\quad
 v_2\det(1\pm\operatorname{Fr}_{r_i}\mid T)\ \leq c_E.
\end{equation}
The constant $c_E$ is independent of the list and of $K$.
\end{proposition}

\begin{proof}
We give the compatibility argument before imposing the Tate condition.
Choose prime ideals $\mathfrak p_1,\ldots,\mathfrak p_s$ with pairwise
disjoint conjugate pairs, including the oriented prescribed primes and
the primes above $2N$, and enough additional split prime ideals to
generate $\operatorname{Cl}(K)$.  Such additional ideals may avoid any
fixed finite set.  Put weight $w_j=1$ at a prescribed odd prime and
$w_j=0$ at every additional prime and every prime above $2N$.
The conjugate weight is $-w_j$.  Let
\[
 R=\ker\bigl(\mathbf Z^s\longrightarrow\operatorname{Cl}(K)\bigr),
\]
choose a basis $m^{(a)}$ of $R$, and choose $\beta_a\in K^\times$ with
$(\beta_a)=\prod_j\mathfrak p_j^{m_j^{(a)}}$.

Work first over $F_\infty=K(\mu_{2^\infty})$ and adjoin $2^n$-th roots
of the finitely many $\beta_a/\overline\beta_a$.  Prescribe on these
roots the multipliers
\begin{equation}\label{rl:kummer-prescription}
 \zeta_{2^n}^{\,-\sum_j m_j^{(a)}w_j}.
\end{equation}
The minus sign fixes the choice of connecting map in the ring-class
sequence and could be reversed throughout.  These prescriptions respect
every Kummer relation.  Indeed, if
\[
 \prod_a(\beta_a/\overline\beta_a)^{b_a}
                         \in(F_\infty^\times)^{2^n},
\]
take its valuation at a prime above each $\mathfrak p_j$ of nonzero
weight.  These odd primes are unramified in $F_\infty/K$.  Thus
$2^n$ divides $\sum_a b_am_j^{(a)}$ for each such $j$, proving the
required compatibility of~\eqref{rl:kummer-prescription}.  Primes over
$2$ have weight zero and impose no division by a ramification index.
Units impose no additional weighted relation.  Kummer duality supplies
the prescribed automorphism of the radical extension over $F_\infty$.

At a rational prime $r$ split in a sufficiently large finite cyclotomic
subfield, the power-residue character of $\mathbf F_r^\times$ is
surjective onto $\mu_{2^n}$.  The split ring-class exact sequence is
\[
 1\longrightarrow
  (\mathcal O_K/r\mathcal O_K)^\times/(\mathbf Z/r\mathbf Z)^\times
 \longrightarrow\operatorname{Pic}(\mathbf Z+r\mathcal O_K)
 \longrightarrow\operatorname{Cl}(K)\longrightarrow1.
\]
The excluded extra-unit fields are the only source of a unit
correction.  The first group is $\mathbf F_r^\times$, with a principal
relation measured by $\beta/\overline\beta$ at the oriented prime over
$r$.  Consequently~\eqref{rl:kummer-prescription} says exactly that its
primitive power-residue character extends to the class group with the
specified values $w_j$.  More explicitly, the inertia group and the
chosen ideal classes present the ring-class group, and the basis of $R$
lists its additional relations.  We have imposed those relations on the
inertia character itself; we have not divided by a class-group order
to extend a previously fixed character.  Thus this construction introduces
no denominator, even when the class number has a large $2$-part.  The
class group and the number of prescribed primes change the finite
Chebotarev extension, but not the integrality of the character.
Since the ring-class extension is unramified at $2N$, value zero on
Frobenius means that the resulting local character is trivial there.
Conjugation acts by inversion.

We now show that imposing these finitely many radical conditions cannot
force either forbidden Tate trace.  In the non-CM case the image of
$G_{K(\mu_{2^\infty})}$ is an open subgroup of
$\operatorname{SL}_2(\mathbf Z_2)$, by the fixed-prime open-image theorem
and the cyclotomic determinant identity
\cite[Chapter~IV, \S2.2, theorem, p.~IV-11, and Chapter~I, \S1.2,
Example~2, p.~I-4]{Serre1968}: the image is the intersection of the
open image of $G_K$ with the determinant-one subgroup.
Its index in the corresponding fixed-$E$ image changes
by at most two when $K$ changes.  Its commutator subgroup contains a
fixed open subgroup $U_E$ of $\operatorname{SL}_2(\mathbf Z_2)$.
One may see uniformity by starting with a principal congruence subgroup
in the fixed image: every index-two subgroup contains its squares, and
the commutators of the resulting fixed congruence subgroup are open.
Because the radical extension over $F_\infty$ is abelian, these
commutators act trivially on it.  Hence every permitted Tate coset
contains a translate of $U_E$.

For completeness, in the CM case let $M$ be the CM field.  The image over
$KM(\mu_{2^\infty})$ contains a fixed open subgroup of the norm-one
Cartan; the index changes by a bounded factor, as $[KM:\mathbf Q]\leq4$.
Choose $\gamma\in G_{KM}$ with
$\chi_{\rm cyc}(\gamma)=5^4=625$.  This is possible uniformly: the
cyclotomic image of $G_{KM}$ has index at most four in
$\mathbf Z_2^\times$, so contains fourth powers.  On the radical Galois
group, conjugation by $\gamma$ is raising to $625$, because all the
radicands lie in $K$.  On the Tate image over $KM$, conjugation is
trivial.  Therefore
\[
 \gamma u\gamma^{-1}u^{-625},\qquad
 u\in G_{KM(\mu_{2^\infty})},
\]
fixes every prescribed radical and acts on $T$ as $\rho_T(u)^{-624}$.
Raising to the fixed nonzero power $624$ carries an open subgroup of the
one-dimensional norm-one Cartan onto an open subgroup.  This proves the
required uniform kernel statement also in the CM case.  The open Cartan
image used here is the CM part of the usual $\ell$-adic image theorem
\cite[Chapter~IV, \S2.2, Remark, p.~IV-13]{Serre1968}.

On a coset of an open subgroup of $\operatorname{SL}_2$, trace is not
identically $2$ or $-2$.  The same holds on a coset of an open norm-one
Cartan subgroup; on the other Cartan-normalizer component trace is zero.
The relevant determinant-one sets are compact and have finitely many
cosets modulo the fixed open subgroup.  Choose in each coset an element
whose trace is neither forbidden value.  The largest of the finitely
many resulting valuations is a constant depending only on $E$.

Finally restrict the prescribed radical automorphism and such a Tate
lift to a finite Galois extension, including as many roots of unity as
required.  Chebotarev gives infinitely many rational primes with this
Frobenius class; choose the place over the prime to realize the chosen
representative.  They split in $K$ and have arbitrarily high prescribed
congruence to $1$.  Since
$\det(1\pm\operatorname{Fr}_r)=1\pm a_r+r$, a congruence precision
larger than the fixed trace bound gives all of~\eqref{eq:h:r-avoidance}.
All these conditions concern a finite extension at each stage.
\end{proof}

Choose $n_i$ large enough that $1+t$ has order dividing $2^{n_i}$ in
$A_i^\times$, and put $\Psi_i=(1+t)^{\xi_i}$.  For example $n_i\geq2i$
is sufficient. Orient the inertia generator to act by $1+t$;
the prescribed split Frobenius exponents remain odd in the limit.
Use global cochains unramified outside $2N$, $r_i$, the active conductor
primes and the derivative primes already retained. At $2N$ use the
Kummer conditions above and elsewhere outside the conductors use the
unramified condition. Reduction modulo $t$, with the unramified
condition restored at $r_i$, gives the ordinary fixed-field complex;
retain its comparison map at every finite stage.

\begin{lemma}[Acyclic relaxed local complexes]\label{rl:local-acyclic}
\label{lem:h:relaxed-local}
For a primitive character constituent, let $C$ denote the limiting
Selmer complex with full conditions at $r$ and at its active good
conductor primes, and the prescribed Kummer, unramified or derivative
conditions elsewhere.  The limiting full
local complexes are acyclic over every height-one residue field of
$\Lambda$ of characteristic zero.  At $(2)$, with residue field
$\Omega=\mathbf F_2((t))$, the $r$-complex is acyclic, as are the
complexes at active split primes with odd prescribed Frobenius exponent
of the deformation character and
at active primes $q$ with $t_q=0$. At each of these acyclic relaxed
places, the full and orthogonal strict conditions agree on the
indicated test. If every relaxed place is acyclic at that test and
the remaining local conditions are self-dual, global duality makes
the Selmer complex self-dual. The Selmer complex
has no global invariants on these tests.
\end{lemma}

\begin{proof}
Use tame inertia and unramified Frobenius, or the Hochschild--Serre
complex for these two procyclic directions.  At $r$, inertia minus one
is multiplication by $t$ times a unit.  It is invertible on all the
indicated fields except the characteristic-zero fiber $t=0$.  There,
$r\to1$, and the two remaining Frobenius differentials have nonzero
determinant by~\eqref{eq:h:r-avoidance}, including either sign of the
quadratic twist.  The bound also ensures that the integral local
torsion at this central fiber has bounded length.

At an active odd conductor prime, inertia of the quadratic character
acts by $-1$.  Its differential is thus $-2$ and is invertible in
characteristic zero.  In characteristic two the quadratic character
disappears.  At a split prime the Frobenius action is $zB$, where
$B\in\operatorname{GL}_2(\mathbf F_2)$ and
$z=(1+t)^a$ has odd $a\in\mathbf Z_2$.  Its coefficient of $t$ is one,
so $z$ is nonconstant.  No nonzero polynomial over the finite constant
field vanishes at $z$.  Thus both Frobenius complexes, with operators
$zB-1$ and $q^{-1}zB-1$, are invertible over $\Omega$; here $q=1$ in
the residue field.  No irreducibility assumption is used.

If $t_q=0$, the rational residual Frobenius $B$ has order three.
At a split place the operator is again $zB-1$, invertible whether
$z$ is nonconstant or $z=1$. At an inert place the ring-class
Frobenius is trivial, since the ideal $(q)$ is principal, and
the Tate Frobenius is $B^2$, which also has no eigenvalue one.
These observations give the additional $t_q=0$ assertion.

The orthogonal full and strict conditions agree at each place where
full local cochains are acyclic. When this holds at every relaxed
place and the remaining local conditions are self-dual,
Equation~\eqref{rl:global-duality} gives global self-duality.
Global invariants in characteristic two are excluded by the nontrivial
$r$-inertia scalar.  In characteristic zero the same argument works
unless the scalar is trivial; in that case $T\otimes\mathbf Q_2$ over
$G_K$ is absolutely irreducible.  For non-CM curves this follows from
the open image, and for CM curves from the open Cartan and its
nontrivial normalizer component, since $K$ is not the CM field.
Lemma~\ref{rl:evaluation} below realizes these actions on the limiting
complex as well. On the self-dual tests, duality pairs degrees
one and two.
\end{proof}

\subsection{The Heegner determinant and its central value}\label{sec:h:specialization}

The point index is not itself the quantity transferred across a growing
family. Each new conductor prime changes the local Selmer terms as well
as the point. We combine these changes in a determinant coordinate.
Its central valuation will be twice the point index, less the finite
Selmer length and the conductor contribution. This formula reduces the
index estimates to lower bounds for the determinant coordinate.

Retain a primitive genus or ring instance from Section~\ref{sec:transfer-criteria},
with point $P_c(e)$ over $K$. Use the conductor characters just
constructed, taking $n_i\ge2i$, and the representations $T\chi_e\Psi_i$.
Retain the same auxiliary primes and the same fixed inner ultrafilter,
now denoted $\mathcal U_{\mathrm{in}}$.

The Selmer complex $C$ is a coefficientwise limit of the corresponding
finite free cochain models along $\mathcal U_{\mathrm{in}}$. It has the integral Kummer condition at
$2N$, the full condition at $r_i$ and the active good conductor primes,
and the unramified condition at all other permitted finite places.
Here the full condition means that no restriction is imposed on local
cochains. Section~\ref{sec:h:finite-models} and Lemma~\ref{rl:local-acyclic}
constructed the local models, their duality maps, and the comparison
at $t=0$. We use these concrete complexes and maps;
cohomologically equivalent complexes without the comparison maps would
not suffice for the later prime choices. Over
$\operatorname{Frac}(\Lambda)$, the modified full local complexes
are acyclic, so the full conditions agree with their orthogonal strict
conditions. Global duality, followed by conjugate transport and the
Weil pairing, therefore gives a perfect pairing
\[
 H^1(C\otimes_\Lambda\operatorname{Frac}(\Lambda))\times H^2(C\otimes_\Lambda\operatorname{Frac}(\Lambda))
 \longrightarrow\operatorname{Frac}(\Lambda).
\]
Here conjugation in $K/\Q$ inverts the ring-class character. These are
the pairing and identification supplied by
Equation~\eqref{rl:global-duality} and Lemma~\ref{rl:local-acyclic}.

Fix the single integer multiplier $L$ of
Lemma~\ref{lem:h:derivative-classes}. It depends on the curve,
parametrization and fixed local types, and clears descent and local
Kummer ambiguities for the complete conductor. In particular its
valuation does not grow with the number of primes. The following
construction and formulas include that multiplier throughout.

At each finite stage transfer the Kummer class of $y_{cr_i}$ from
$H_{cr_i}$ to $K$ with the specified scalar coefficients, using
$Ly_{cr_i}$. Complete the finite coset sum before contracting its
target cocycle and local Kummer lifts into the models over $K$.
The resulting cycles $P_i$, with their contraction witnesses, give
$P\in H^1(C)$ by Lemma~\ref{rl:finite-limits}.
Local Kummer conditions at $2N$ are preserved because $\Psi_i$ is
trivial there. At the center of the $i$th finite model, the norm relation
has the factor
\begin{equation}\label{eq:h:trace-center}
 \begin{aligned}
 b_i&=a_{r_i}-\chi_e(\Fr_{\mathfrak r_i})
            -\chi_e(\Fr_{\bar{\mathfrak r}_i}),
 \qquad r_i=\mathfrak r_i\bar{\mathfrak r}_i,\\
 P_i\big|_{t=0}&=\varepsilon_i b_i\kappa(LP_c(e)).
 \end{aligned}
\end{equation}
where $\kappa$ denotes the Kummer class in the central comparison
and $\varepsilon_i$ is a unit; the equality is taken modulo the finite
coefficient precision. Since $\chi_e$ is defined over $\Q$, its values
at the conjugate Frobenius elements agree; explicitly,
$b_i=a_{r_i}-2\chi_e(r_i)$.
To retain the central identity at the cochain level, let $B_i$ be the
reduction of a fixed finite model of the ordinary Kummer complex over
$K$, and let $\beta_i:B_i\to C_i/(t)$ be its comparison map to the
relaxed central complex. Choose a cycle $q_i$ representing
$\kappa(LP_c(e))$. The finite norm relation gives a boundary witness
$v_i$ in the target model such that
\[
 P_i\bmod t-\varepsilon_i b_i\beta_i(q_i)=d v_i.
\]
Retain $\beta_i,q_i,v_i$, and the target contraction homotopies on
the same stages and at the same coefficient precision as $P_i$.
Retain these units in the same ultrafilter
$\mathcal U_{\mathrm{in}}$ as all the comparison maps. Then
$b=\lim_{\mathcal U_{\mathrm{in}}}b_i$ has bounded valuation by
Equation~\eqref{eq:h:r-avoidance}, and
$\varepsilon=\lim_{\mathcal U_{\mathrm{in}}}\varepsilon_i$ is a unit. Thus
$P|_{t=0}=\varepsilon b\kappa(LP_c(e))$, with $b\ne0$.
The special complex differs from the ordinary Kummer complex only by
the active-prime relaxed quotients and the limit of the singular
quotient at the single moving prime $r_i$, the latter of bounded length.
Inflation from the complexes omitting $r_i$ gives this comparison. All modified
local complexes are rationally acyclic at the center, so pairings agree
there as well.

If the product has a simple zero, its rational special fiber has
one-dimensional cohomology in degrees one and two, and $P$ specializes
nontrivially by Gross--Zagier. The generic dimensions are at most these
special dimensions. Moreover, the long exact sequence for multiplication
by $t$ shows that $H^1(C[1/2])[t]=0$, since the special $H^0$ vanishes.
A torsion class that disappeared over the generic field would be killed
by a power of $t$ after localization at $(t)$; hence the nonzero special
class cannot disappear. Thus the generic ranks are one in degrees one
and two, $P$ spans $H^1$, and the rational rank is constant at $t=0$.
If only base Selmer corank one is known, the same dimension bound gives
generic ranks at most one; a zero generic point is allowed.

\begin{definition}\label{def:h:determinant-coordinate}
Choose an integral basis of $\det^{-1}_{\Lambda}C$. In generic rank one,
the element formed by $P$ in $H^1$ and the functional
$\langle P,\cdot\rangle$ on $H^2$ determines a scalar
$u(C,P)\in\operatorname{Frac}(\Lambda)$ in that basis. Put $u=0$ if
$P=0$ or the generic rank is zero. Integral changes of determinant basis
multiply $u$ by a unit.
\end{definition}

In the simple-zero case, constancy of the rational cohomology rank at
$t=0$ makes $u$ regular in the local ring $\Lambda_{(t)}$. Its value
$u(0)\in\Qtwo$ is therefore already defined. This specialization takes
place at $(t)$; the stronger assertion $u\in\Lambda[1/2]$ will follow
from the divisor estimates below.

The determinant has a simple integral interpretation. On a rank-one
free Selmer line, the point and its paired functional each contribute
the point's index. Torsion in degree two subtracts its length from the
coordinate in the inverse determinant lattice. For example, the complex
$[\Ztwo\xrightarrow{2^a}\Ztwo]$ in degrees one and two has inverse
determinant lattice mapping to $2^a\Ztwo$ in
the rational determinant. The coordinate of a fixed rational tensor
therefore loses $a$ in valuation. The next lemma carries out this
calculation with the actual local Kummer conditions. It also compares
the Selmer length over $K$ with those of the two rational twists.

\begin{lemma}[Central determinant valuation]
\label{lem:h:central-valuation}
For a simple-zero instance,
\begin{equation}\label{eq:h:central-determinant}
 \begin{aligned}
 v(u(0))&=2j(P_c(e))
   -\length\Sha(E^{(e)}/K)[2^\infty]
   -\sum_{q\mid c}\sum_{v\mid q}t_v+O_E(1),\\
 t_v&=\dim_{\F_2}E[2]^{\Fr_v}.
 \end{aligned}
\end{equation}
Furthermore,
\begin{equation}\label{eq:h:weil-restriction-length}
 \length\Sha(E^{(e)}/K)[2^\infty]
       +\sum_{q\mid c}\sum_{v\mid q}t_v
 =s(e)+s(ek)+\sum_{q\mid k}t_q
                         +2\sum_{q\mid c}t_q+O_E(1).
\end{equation}
The constants include the one fixed class multiplier and the bounded
local contribution at the moving prime $r_i$.
\end{lemma}

\begin{proof}
The ordinary integral Kummer complex over $K$ is self-dual. Its $H^1$ is
the compact Selmer module, and the torsion in $H^2$ has the length of the
finite quotient of the discrete Selmer group by its divisible subgroup.
This follows by derived base change to $\Qtwo/\Ztwo$: the local
complexes give exactly the discrete Kummer conditions. In the present
simple-zero case that finite quotient is
$\Sha(E^{(e)}/K)[2^\infty]$. Duality pairs the free parts in degrees
one and two perfectly. Torsion in degrees one and three costs a bounded
amount, by global torsion and duality. The uniform torsion bound here comes
from a fixed nonidentity Tate scalar acting trivially on every ring-class
field in the construction: its scalar minus one annihilates that torsion.
Lemma~\ref{rl:kernel} constructs this scalar, and the descent argument
of Lemma~\ref{lem:h:derivative-classes} uses the same bound.

At an active odd conductor place inertia acts by minus identity.
All integral $H^1$ is Kummer, and the relaxed quotient contributes only
$H^2(K_v,T^e)$ in degree two. Its length is $t_v$, by local duality
and the fixed $2$-torsion. At the moving prime $r_i$ the contribution
is bounded by
Equation~\eqref{eq:h:r-avoidance}. The generic pairings specialize to
the ordinary Kummer pairing because the modified local complexes are
rationally acyclic at the center. Thus the free determinant tensor has
valuation twice the point index, whereas torsion in degree two
contributes negative length to the inverse determinant. Together with
Equation~\eqref{eq:h:trace-center}, this proves
Equation~\eqref{eq:h:central-determinant}. The specialization is at
constant rational cohomology rank, so no specialization torsion is
being substituted for the point.

For Equation~\eqref{eq:h:weil-restriction-length}, apply the isogeny
invariance of the BSD quotient to quadratic Weil restriction and
$E^{(e)}\times E^{(ek)}$
\citep[Theorem~2.1 and the preceding discussion]{Tate1966}.
All the groups in this
comparison are finite apart from the known rank-one lattices; no BSD
equality is assumed. The forward and reverse isogenies have composites
multiplication by $2$. Consequently the rank-one lattice indices,
height ratios, torsion ratios and period ratios contribute bounded
$2$-powers. For periods, the pullback indices of N\'eron differentials
at $2$ are bounded by those composites; real kernels and component
ratios are bounded as well.

The remaining terms are the Tamagawa numbers. A ramified quadratic
twist of good reduction at an odd prime $q$ has type $I_0^*$ and
component number one plus the number of roots of its reduced auxiliary
cubic \citep[Theorem~4.1, Table~4 and proof, Case~3]{BarriosEtAl2025}.
To identify this number, write a good integral model as $y^2=f(x)$,
retaining the quadratic term, and the twisting parameter as $d=\pi u$.
The auxiliary cubic is $\bar u^3\bar f(T/\bar u)$; it is separable
and its roots correspond to the nonzero rational two-torsion points.
Thus the component number is $\#E[2]^{\mathrm{Fr}_q}$ and its
$2$-valuation is $t_q$. This argument also applies in residue
characteristic $3$. Its connected reduction is unipotent. Some inertia element acts
by $-1$ on the twisted $2$-adic representation, so its local
$2$-primary torsion is precisely its Frobenius-fixed $2$-torsion.
Over an unramified extension the same formula holds with the
Frobenius of the enlarged residue field. At unramified
or split places, Weil-restriction component numbers are the products
over $v\mid q$. At a prime $q\mid k$ outside $2N$, the curve over
$K$ has good reduction and its Weil restriction has connected special
fiber. Indeed, Weil restriction of its N\'eron model is the N\'eron
model of the restriction of the curve
\citep[Section~7.6, Proposition~6]{BoschLutkebohmertRaynaud1990}.
Its special fiber has the original connected elliptic fiber as quotient,
with connected vector-group layers from ramification. Exactly one rational twist
has the ramified component factor. At a conductor prime $q\mid c$,
the two rational twists give a combined valuation $2t_q$, and the
corresponding contribution over $K$ is $\sum_{v\mid q}t_v$. All
places above $2N$ have bounded local contributions. These are precisely
the terms of Equation~\eqref{eq:h:weil-restriction-length}.
\end{proof}


In a simple-zero genus instance, the fixed partner absorbs the terms involving $h_*$,
while $c=1$ and $k=hh_*$. The displayed formulas therefore identify
\[
 v(u(0))=2j(P_1(h_*))-2w(h)-s(h)+O_{E,h_*}(1).
\]
In a simple-zero ring instance with $\omega(k)\le a$, the discriminant term
is bounded and the conductor term is $4w(h)$, giving
\[
 v(u(0))=2j(P_{|h|}(h))-4w(h)-s(h)-s(hk)+O_{E,a}(1).
\]
The ring identity also applies to the interpolation families in which
an added nonsplit prime has $t_q=0$; such a prime contributes no weight.
The all-split hypothesis is needed for the uniform ring lower bound.
These formulas explain both the lower-bound statements and the
normalization in Heegner interpolation. To prove the lower bounds,
we must show that $u$ has no negative divisor away from $(2)$ and
has a uniformly bounded negative divisor at $(2)$ in the stated
cases. We will construct integral local switches that preserve its determinant
valuation, and use inner detection to choose switches lowering excess
Selmer rank.
The uniform ring estimate will require an additional characteristic-zero
limit; the universal binary family will instead require a bounded
clearing factor after the new local blocks have been removed.

The main intermediate estimate is
\begin{equation}\label{eq:h:divisor-bounds}
 \begin{aligned}
 v_{\mathfrak p}(u)&\ge0
       &&\bigl(\operatorname{ht}\mathfrak p=1,
                         \ \mathfrak p\ne(2)\bigr),\\
 v_{(2)}(u)&\ge-O_{E,h_*}(1)
       &&\text{in genus data with irreducible }E[2],\\
 v_{(2)}(u)&\ge-O_{E,a}(1)
       &&\text{in ring data with }\omega(k)\le a
                         \text{ and all active primes split.}
 \end{aligned}
\end{equation}
To construct the switches used in these estimates, we first retain
enough Galois evaluation data to select their derivative primes.

\subsection{Evaluation and the exact kernel}

The determinant estimates require switches that detect the excess
Selmer dimension at a height-one test. To choose
switching primes, we must still detect a nonzero limiting Selmer class
by Galois evaluation. The next two lemmas first retain that detection
under coefficient specialization, then restrict it to sequences acting
trivially on the full Tate module, all ring-class fields, and the
quadratic characters in use.

\begin{lemma}[Evaluation on arbitrary Galois sequences]\label{rl:evaluation}
\label{lem:h:evaluation}
Fix $K$ and a sequence of the finite cochain constructions above.  Keep
the contraction maps into the original global cochains.  Let
$\Gamma_K=\prod_i G_K$ and
$\Gamma=\Gamma_K\rtimes\langle\tau\rangle\subset\prod_iG_{\mathbf Q}$,
with simultaneous complex conjugation $\tau$.  Thus $\Gamma_K$ has
index two in $\Gamma$; the larger unrestricted product of rational
Galois groups is not used for induction.  The matrix limit defines a
coefficient action of $\Gamma_K$ on $\Lambda^2$.  For every field $L$
receiving a map from $\Lambda$, there is a natural injection
\[
 H^1(C\otimes_\Lambda L)\longrightarrow
 H^1_{\rm abs}(\Gamma_K,L^2),
\]
where the right side denotes abstract crossed-homomorphism classes.
Here $C$ may be the global complex or a Selmer complex whose local
conditions contain all local degree-zero invariants and inject into
local $H^1$.  The same statement holds for the induced representation
of $\Gamma$.

The evaluation maps and their degree-two cochain identities exist
for every group sequence and commute with coefficient extension.
\end{lemma}

\begin{proof}
Let $C_i$ be a finite model and let $\iota_i$ be its chosen map to the
original global cochains, using the projection from the Selmer cone
when appropriate.  For every $g=(g_i)$ and $h=(h_i)$ define matrices
\[
 e^0_i:C_i^0\to A_i^2,\quad
 e^1_{g,i}:C_i^1\to A_i^2,\quad
 e^2_{g,h,i}:C_i^2\to A_i^2
\]
by evaluating $\iota_i$ in degrees zero, one, and two.  Their source
ranks are fixed and their entries belong to $A_i$.  Thus each has a
coefficientwise ultrafilter limit, for every choice of $g,h$.
The same fixed ultrafilter defines these limits for all $g,h$.  Likewise
$\rho(g)=\lim_i\rho_i(g_i)$ is a representation.  The original
cochain identities imply
\begin{align}
 e^1_g d&=(\rho(g)-1)e^0,\label{rl:eval-boundary}\\
 e^2_{g,h}d&=\rho(g)e^1_h-e^1_{gh}+e^1_g.
                         \label{rl:eval-cocycle}
\end{align}
All identities pass to the limit because their matrix expressions
have finitely many terms.  A cycle $x\in C^1\otimes L$ consequently
gives the crossed homomorphism $g\mapsto e^1_g(x)$, and a boundary
gives a coboundary.

We prove injectivity in a way that also controls base change.  At each
finite stage let $k=\mathbf F_2$ be the residual field.  Choose a
space of representatives for residual $H^1$.  Evaluate pairs $(c,m)$,
where $c$ is in this space and $m\in k^2$, by
$b_c(g)-(\rho(g)-1)m$.  The intersection of the kernels of all
these evaluations consists precisely of pairs $(0,m)$ with $m$ an
invariant.  Since the vector space is finite dimensional, a list of
at most $\dim H^1+2$ elements already has this intersection.  For
these elements evaluation is a chain map
\[
 C_i\longrightarrow D_i=
 [\,A_i^2\xrightarrow{((\rho_i(g_j)-1))_j}(A_i^2)^d\,]
\]
in degrees zero and one, inducing an isomorphism on residual $H^0$
and an injection on residual $H^1$.  For Selmer complexes this uses
the exact sequence of~\eqref{rl:selmer-definition}: the indicated
local degree-zero isomorphisms give $H^0(C_i)=H^0_{\rm glob}$ and
$H^1(C_i)\hookrightarrow H^1_{\rm glob}$ residually.

Retain these matrices and their chosen group-element sequences.
The residual global and local complexes, and the local conditions,
have no negative cohomology.  The cone definition therefore gives
$H^n(C\otimes k)=0$ for $n<0$.  Put $Q=\operatorname{Cone}(C\to D)$.
The long exact sequence gives $H^n(Q\otimes k)=0$ for $n<-1$;
its groups in degrees $-1$ and $0$ vanish because the detecting map is
an isomorphism on $H^0$ and injective on $H^1$.  Any $H^3(C\otimes k)$
contributes in degree two of this cone, so causes no negative-degree
obstruction.  Successively cancel the unit differential blocks of
$Q$.  This gives a homotopy-equivalent bounded free complex starting
in degree one.  For any $\Lambda$-algebra $L$, its ordinary tensor
product computes $Q\otimes^{\mathbf L}_\Lambda L$ and still starts
in degree one; no flatness is needed.  Thus the long exact sequence
makes $H^1(C\otimes L)\to H^1(D\otimes L)$ injective.  A crossed cocycle that is a coboundary
on all of $\Gamma_K$ is in particular a coboundary on the detecting
list, proving the desired injection.  For induction the same proof
has coefficient rank four and the detecting bound is $\dim H^1+4$.
\end{proof}

\begin{lemma}[The exact common kernel]\label{rl:kernel}
In the notation of Lemma~\ref{rl:evaluation}, let $H$ be the group of
sequences each of whose entries acts trivially on the full Tate module,
on all ring-class fields of $K$, and on the quadratic characters in
use.  Restriction from the crossed classes in that Lemma to $H$ is
injective at every characteristic-zero test.  It is also injective
at the residual $(2)$ test.  The characteristic-zero annihilator
needed for this assertion can be chosen with a fixed finite
$2$-valuation depending only on $E$, also as $K$ varies.
\end{lemma}

\begin{proof}
The Tate image contains an open subgroup of scalar homotheties, both
in the CM and non-CM cases.  Choose once and for all
$g\in G_{\mathbf Q}$ with Tate action $u\,1$, where
$u\in1+4\mathbf Z_2$ and $u\ne1$.  Its square belongs to every
quadratic subgroup.  Put
\[
 z=g^2\tau g^2\tau^{-1}.
\]
The generalized dihedral property of every ring-class extension
implies that $z$ is ring-class trivial.  It is also trivial on every
quadratic character, and its Tate action is the fixed scalar $u^4$.
For every sequence in $\Gamma_K$, its commutator with the constant
sequence $z$ is in the exact group $H$, since its Tate commutator and
all its abelian character images are trivial.  Thus $z$ is central in
$\Gamma_K/H$ and acts on the coefficients by $u^4$.

Here the elementary central-scalar argument suffices.  If a crossed
cocycle $b$ is defined on a group with such a central element, comparison
of $b(zg)$ and $b(gz)$ gives
\[
 (u^4-1)b(g)=(g-1)b(z).
\]
In characteristic zero this makes $b$ a coboundary.  Inflation--restriction
then gives injectivity on $H$.  The valuation of $u^4-1$ is fixed.

In characteristic two choose, at each stage, inertia at $r_i$ mapping
onto the chosen generator of $\xi_i$. Its Tate action is trivial and
its ring-class images commute, so the resulting sequence is central
in $\Gamma_K/H$ and acts on $M$ by $1+t$. Its scalar minus one is
invertible in $\mathbf F_2((t))$.

For the induced representation $W$, first restrict to $\Gamma_K$.
There $W=M\oplus M^\tau$, and this inertia sequence acts on the
summands by $1+t$ and $(1+t)^{-1}$. Hence $W^{\Gamma_K}=0$.
Inflation--restriction for the index-two inclusion gives an injection
$H^1(\Gamma,W)\to H^1(\Gamma_K,W)$, even in characteristic two.
Apply the central-scalar calculation separately to the two summands
to obtain injection on restriction to $H$. The same argument works
in characteristic zero with the fixed scalar $u^4$.
Throughout, $H$ is defined by exact actions at every stage.
\end{proof}

\begin{lemma}[Irreducibility on the characteristic-zero tests]
\label{rl:irreducibility}
On an inner characteristic-zero test where $\Psi^2\ne1$,
the four-dimensional representation induced from $T\chi\Psi$ is
absolutely irreducible.  If $\Psi^2=1$ on an inner test, it splits
as two inequivalent absolutely irreducible quadratic twists.
\end{lemma}

\begin{proof}
The two induced summands are distinguished by $r$-inertia, whose
scalars are $\Psi$ and $\Psi^{-1}$.  On either summand, multiplying
the Tate matrices by scalars does not alter invariant subspaces.
The Tate representation restricted to $G_K$ is absolutely
irreducible under the exclusions already specified.  These facts
give absolute irreducibility of the induction when the two inertia
characters differ.

If $\Psi^2=1$, extension across conjugation gives the two quadratic
twists.  They are inequivalent, since a quadratic self-twist of the
Tate representation here would make its restriction to $G_K$
reducible over an algebraic closure.  In the CM case this would
force $K$ to be the CM field, which was excluded.
\end{proof}

\begin{lemma}[Simultaneous finite Chebotarev realization]
\label{rl:chebotarev-evaluation}
For any element $u$ of an exact kernel in Lemma~\ref{rl:kernel},
derivative primes can be chosen, at increasing finite precision,
whose rational Frobenius is $\tau u$ on all retained data.  Their
finite localization is evaluation on $(\tau u)^2$.  The assertion
is compatible with any fixed finite list of previous derivative
operations.
\end{lemma}

\begin{proof}
At one Artin stage, a coordinate of a cochain in the image of a
finite model is a locally constant function.  The finitely many
coordinates being used factor through one finite quotient of the
absolute Galois group.  Enlarge it to contain the Tate action at
the required precision, the finitely many ring-class fields and
quadratic characters currently in use, and all old finite data.
Chebotarev realizes the conjugacy class of $\tau u$ on this finite
quotient, with a chosen place realizing the desired representative
\citep[Chapter~V, Theorem~3.23]{MilneCFT2020}.
It gives a rational prime inert in $K$, outside the old finite set.
Because $u$ is exactly Tate-trivial, its Frobenius tends to
$\rho_T(\tau)$; in particular $a_\ell\to0$ and $\ell\to-1$.
Because $u$ is exactly ring-class trivial, the square Frobenius on
the $K$-prime is the identity on every retained ring-class field.
The unramified local coordinate is therefore the stated
Frobenius-square evaluation.

Repeat at each finite stage with increasing precision.  All
conditions remain finite, and no density bound or prime-size
bound is required.
Previously imposed finite diagrams are unchanged.
\end{proof}

In induced-module notation, restriction to the exact kernel is additive.
For a crossed cocycle $b$ and an exact-kernel element $u$,
\[
 b((\tau u)^2)=(1+\tau)b(u).
\]
Indeed $\tau^2=1$, $u$ acts trivially on the coefficients, and the
cocycle identity cancels the two occurrences of $b(\tau)$.  Projection
to either induced summand identifies $\operatorname{im}(1+\tau)$
with a two-dimensional coefficient module, also in characteristic two.

\subsection{The local switch and its integral determinant}
\label{rd:local-subsection}

Fix a complex conjugation $\tau$ and write $C_\tau$ for its action on $T$.
The next two Lemmas construct the integral local switches and the classes
that pass through them. Their finite-level formulas are needed before
applying the limit constructions above.

\begin{lemma}[The split local model]\label{rd:local-model}\label{lem:h:switching-pairing}
Let a sequence of inert good primes \(\ell\) have rational Frobenius
lifts tending on \(T\) to a fixed complex conjugation \(C_\tau\).
At the unique prime of \(K\) above \(\ell\), the limiting local complex
has
\begin{equation}
 H^0=M,\qquad H^1=F\oplus S=M\oplus M,\qquad H^2=M,
 \label{rd:local-cohomology}
\end{equation}
with zero differentials in a free split model.  Here \(M=T\otimes\Lambda\)
with the relevant scalar coefficient structure.  The two degree-one
coordinates of a cocycle are
\begin{equation}
 f(c)=c(\gamma^2),\qquad s(c)=c(\sigma),
 \label{eq:h:local-coordinates}
\end{equation}
where \(\sigma\) is an oriented pro-\(2\) tame generator.  The plane
conditions \(F=\{s=0\}\) and \(S=\{f=0\}\) include \(H^0\) and exclude
\(H^2\).  Each is its own orthogonal complement for conjugate-Weil
duality.  The strict condition \(H^0\) and the relaxed condition
\(H^0\oplus H^1[-1]\) are mutually dual.  All four inclusions can be
equipped with compatible integral duality triangles.
\end{lemma}

\begin{proof}
The prime-to-\(\ell\) local cochains are computed by the tame inertia and
residue-field directions.  In these two directions the differentials
are formed from inertia minus one and Frobenius minus one, with the
usual Frobenius action on tame inertia.  On \(M\), inertia is trivial at
the derivative prime.  The square of a rational Frobenius is trivial
on every ring-class character.  In the limit it is also the identity on
\(T\), while \(\ell^2\) tends to one.  Thus both differentials become zero,
giving \eqref{rd:local-cohomology}.

Here is a finite-precision check of the pairing, needed at \(2\).
At a retained quotient killed by \(2^a\), impose also
\(q=\ell^2\equiv1\pmod{2^{a+1}}\), as well as trivial
Frobenius-square action on the coefficient module at this quotient.
The tame two-generator complex then has zero differentials,
with the coordinates in \eqref{eq:h:local-coordinates}.  Its degree-one
scalar generators are the unramified and tame characters.
More explicitly, use the presentation
\(\phi\sigma\phi^{-1}\sigma^{-q}=1\).
After augmentation with trivial coefficient action its Fox
resolution has terms of ranks \(1,2,1\), with
\(d^0=0\) and \(d^1=(0,1-q)\), in the ordered finite and
tame coordinates.  The diagonal from the same resolution
has zero pure-finite term, unit mixed terms, and pure-tame
coefficient \(\pm\binom q2\).  The latter coefficient is
the augmented second derivative of \(\sigma^{-q}\)
together with its adjacent \(\sigma\) in the relation.
Both \(1-q\) and \(\binom q2\) vanish modulo \(2^a\)
under the imposed congruence.  Thus this is an actual
split complex and an actual pure-plane zero pairing
at the retained precision.
The field contains the \(2^{a+1}\)-st roots of unity, so
\(-1\) is a \(2^a\)-th power.  The norm-residue identity
\((x,x)=(x,-1)\) for the \(2^a\)-Hilbert symbol therefore
makes the square of each scalar generator zero.
The two generators form a basis; local Tate duality makes
their cross pairing a unit.  The same calculation applies
coefficientwise to the free Artin coefficient module, including
its nilpotent \(t\) and quadratic variables.  Thus the pure
plane restrictions are already zero at every retained quotient.
All these congruences can be imposed simultaneously with the
derivative Frobenius conditions by increasing their precision.
Taking the compatible limit proves integral vanishing of the
pure squares.  It is not deduced from skew symmetry in
characteristic two.

Equivalently, in the limiting tame presentation
\(\phi\sigma\phi^{-1}=\sigma^q\) with \(q\to1\), the
rank-\((1,2,1)\) resolution specializes to the exterior
resolution for two procyclic directions.  The finite
Hilbert-symbol calculation specifies the cup product on
its split cohomology and rules out a hidden square term.
Conjugation transports the finite Tate coordinate by \(C_\tau\) and the
singular coordinate by \(-C_\tau\), with possible scalar units coming from
the orientations.  Consequently the cross pairing is
\begin{equation}
 \beta(x,y)=\epsilon\,\langle x,C_\tau y\rangle,\qquad
 \epsilon\in\Lambda^\times,\quad x\in F,\ y\in S,
 \label{eq:h:cross-pairing}
\end{equation}
where \(\langle\, ,\,\rangle\) is the perfect alternating Weil form.
This proves the assertions about orthogonal complements.  Since the
restrictions of this explicit pairing to either pure plane are zero as
maps of the split complexes, their isotropy homotopies may be chosen to
be zero.  Only a product of two degree-one terms could map to the
degree-two local invariant target; the degree-zero terms introduce
no additional restriction.  These zero homotopies agree on the
strict subcomplex.  These are the actual
orthogonality data required in the Selmer duality construction
\cite[Theorem~6.3.4]{Nekovar2006}; the compatibility assertion is therefore
an assertion about these displayed complexes and maps.  Lemma~\ref{rl:finite-limits} is used to retain this finite diagram together
with the unchanged local data.
\end{proof}

For a finite set $I$ of derivative primes, let $C_I$ be obtained
from $C$ by replacing the finite plane condition $F$ by $S$ at each
prime of $I$, retaining all other local conditions. The classes below
use the same multiplier $L$ already included in $P$; in particular,
$P_{\varnothing}=P$.

\begin{lemma}[Descent and the switch at \(2\)]\label{rd:derivative}\label{lem:h:derivative-classes}
Suppose that the finite ring-class transfers are performed before
retaining their target cycles and witnesses as in
Lemma~\ref{rl:finite-limits}. There is one
integer \(L\ne0\), depending on the fixed curve, parametrization and
fixed local types, such that the classes obtained from the Heegner
derivatives of \(L\) times the point satisfy, for every finite switched
set \(I\),
\begin{equation}
 P_I\in H^1(C_I),\qquad
 f(P_{I\cup\{\ell\}})=0,\qquad
 s(P_{I\cup\{\ell\}})=C_\tau f(P_I).
 \label{eq:h:finite-singular}
\end{equation}
The same \(L\) works for every number of active and derivative primes.
The previously imposed singular conditions are retained.
\end{lemma}

\begin{proof}
We first descend the derivative classes using one multiplier for the
complete conductor, and verify their Kummer conditions at $2N$. We then
compute both local coordinates in Equation~\eqref{eq:h:finite-singular}.
The last step retains these classes and identities in the same bounded
target cochain diagrams used for evaluation.

Choose the orientations coherently.  The relative ring-class group at
an inert prime \(\ell\) is cyclic of order \(\ell+1\).  For distinct new
primes these groups form a product: the ring-class exact sequence has
kernel the product of the local residue-unit quotients, and the only
global units here are \(1,-1\), whose images in these quotients are
trivial.  Set
\[
 D_\ell=\sum_{i=1}^{\ell} i\sigma_\ell^i,\qquad
 D_I=\prod_{\ell\in I}D_\ell.
\]
Direct multiplication gives
\begin{equation}
 (\sigma_\ell-1)D_\ell=(\ell+1)-
             \sum_{i=0}^{\ell}\sigma_\ell^i.
 \label{rd:derivative-identity}
\end{equation}
The Hecke correspondence and reduction of its CM isogenies give, for
every class-group translate,
\begin{equation}
 \operatorname{Tr}(y_{c r n})=a_\ell y_{c r n/\ell},
 \qquad
 \widetilde y_{c r n}=\operatorname{Fr}_\ell
                       \widetilde y_{c r n/\ell}.
 \label{eq:h:heegner-norm-reduction}
\end{equation}
These are the point relations, not an application of a residual
irreducibility theorem.  A fixed multiple of the parametrization removes
cuspidal terms.

There is a uniform bound on the \(2\)-primary torsion over the union of
the ring-class fields in use.  To see this, choose a fixed nonidentity
scalar homothety of infinite order in the \(2\)-adic Tate image.
Square a lift so that it
belongs to every quadratic subgroup under consideration.  Multiply it
by its conjugate under complex conjugation.  Its action on each
ring-class field is trivial, because conjugation inverts the abelian
ring-class action, while its action on \(T\) is still a fixed
nonidentity scalar.  If \(2^e\) annihilates its scalar minus one up to
a unit, then \(2^e\) kills the torsion in question.  The existence of
this fixed scalar, including the CM case, is the Tate-image input
used in Lemma~\ref{rl:kernel}.

Work first with coefficients modulo \(2^a\).  Require \(a_\ell\) and
\(\ell+1\) to vanish to a precision substantially larger than \(a\).
Equations \eqref{rd:derivative-identity} and
\eqref{eq:h:heegner-norm-reduction} then make the derivative Kummer class
invariant under the relative ring-class group.  In the
inflation--restriction sequence, both the obstruction to descent and
the ambiguity of descent are killed by \(2^e\), since their
coefficient module is killed by \(2^e\).  Multiplication by \(2^{2e}\)
therefore gives a descent independent of this ambiguity: descend
\(2^e\) times the class and multiply its lift by \(2^e\).
There is no repetition of this factor for individual primes; descent
is performed once from the field with the complete conductor.

The global descent has cost only the fixed factor $2^{2e}$, regardless
of the number of derivative primes. It remains to ensure that the
descended class satisfies the Kummer condition at every fixed bad place;
this is a separate local descent issue.

At a place over \(2N\) these descent extensions are unramified.
Here is the local descent argument, valid also for even extension degree
and for toric or additive reduction.  Let $F'/F$ be a finite unramified
extension with group $G$, and write $k'/k$ for its residue extension.
The successive formal-group quotients $E_j(F')/E_{j+1}(F')$ are additive
copies of $k'$.  The normal basis theorem makes each a cohomologically
trivial $G$-module, without a restriction on $|G|$.  Every finite
quotient $E_1(F')/E_j(F')$ is therefore cohomologically trivial.
The inverse systems of these finite groups and their invariant groups
satisfy the Mittag--Leffler condition.  Taking inverse limits in finite
$G$-cochains consequently gives $H^1(G,E_1(F'))=0$.
Let $\mathcal E/\mathcal O_F$ be the N\'eron model and put
$\widetilde E^0=\mathcal E_k^0$, a smooth geometrically connected
commutative algebraic group over the finite field $k$.
Lang's finite-extension cocycle theorem
\citep[Proposition~3, p.~559]{Lang1956} gives
$H^1(G,\widetilde E^0(k'))=0$.
The smooth reduction sequence then gives $H^1(G,E^0(F'))=0$.
Let $\Phi=\mathcal E_k/\widetilde E^0$ be the finite \'etale
component-group scheme. Each fiber over a point of $\Phi(k')$
is a torsor under $\widetilde E^0_{k'}$, and hence has a $k'$-point
by \citet[Theorem~2, pp.~557--558]{Lang1956}.
Smooth lifting therefore makes $E(F')\to\Phi(k')$ surjective. Hence
\[
 H^1(G,E(F'))\lhook\joinrel\longrightarrow
 H^1(G,\Phi(k')).
\]
Existence of the smooth finite-type model and smooth lifting follow from
\citet[Section~1.3, Corollary~2, and Section~2.3, Proposition~5]{BoschLutkebohmertRaynaud1990}.
The right side is killed by the exponent of $\Phi(\overline k)$.
By inflation--restriction, the left side is exactly the kernel of
$H^1(F,E)\to H^1(F',E)$.  Thus the same fixed exponent makes a
class in $H^1(F,E[2^a])$ Kummer whenever its restriction is Kummer.
N\'eron models commute with unramified base extension
\citep[Section~7.2, Theorem~1 and Corollary~2]{BoschLutkebohmertRaynaud1990},
so their geometric component groups are unchanged. Only the finitely many fixed local twisting
types at $2N$ occur, so the product of these component exponents is
bounded independently of the conductor and of the extension degree.
Include this product, the single descent factor $2^{2e}$, and the
parametrization factor in \(L\).  Thus the descended class has the
integral Kummer condition at \(2N\).
No good-ordinary or semistability condition was used.  Transfer with
the scalar character preserves this condition because \(\Psi=1\)
there.

The single multiplier $L$ now gives descent with the required Kummer
conditions. We next prove the finite-to-singular identity itself. Both
coordinates matter: the singular coordinate must become $C_\tau f$,
while the new finite coordinate must vanish at the retained precision.

Here is the local calculation, including the division by two.  Apply
the derivatives at all primes other than \(\ell\), and write \(X,Y\)
for the resulting lower and upper points.  Put \(F_\ell=\operatorname{Fr}_\ell\).
The lower field splits at the \(K\)-prime \((\ell)\), since that ideal
is principal.  Reduction therefore gives \(F_\ell^2X=X\).
The following Frobenius-polynomial identities are taken on the reduction;
good odd-prime reduction is injective on the $2$-primary torsion
coordinates being compared.
If \(Q\) is a \(2^a\)-division point of \(LX\), the finite coordinate
of its Kummer class is
\begin{equation}
 f=(F_\ell^2-1)Q=
       \bigl(a_\ell F_\ell-(\ell+1)\bigr)Q.
 \label{rd:finite-formula}
\end{equation}
We give the descent-cocycle calculation of the other coordinate.
Choose \(Z\) with \(2^aZ=LD_\ell Y\), and choose the descended
cocycle \(b\) so that on the upper-field Galois group it is
\(b(g)=(g-1)Z\).  A change by the already bounded descent ambiguity
has no effect on the eventual limit.  The crossed identity gives
\[
 U=(\sigma_\ell-1)Z-b(\sigma_\ell)
                  \in E(H_{crn}).
 \]
For \(g\) in the upper-field Galois group, compare the two
expansions of \(b(g\sigma_\ell)\); they give \(gU=U\).
Moreover,
\[
 2^aU=L\bigl((\ell+1)Y-a_\ell X\bigr).
 \]
Both integers \(\ell+1,a_\ell\) are divisible by \(2^a\).
Consequently
\[
 U_0=L\frac{\ell+1}{2^a}Y-
                      L\frac{a_\ell}{2^a}X
 \]
is an actual upper-field point with the same \(2^a\)-multiple
as \(U\).  Thus \(U-U_0\) is upper-field torsion of bounded
exponent.  This comparison with an actual integer quotient
is what bounds the ambiguity; arbitrary division points
would not give this bound.  Reduction makes
\(\sigma_\ell\) act trivially and
\(\widetilde Y=F_\ell\widetilde X\).  Therefore
reduction of \(b(\sigma_\ell)=-U\) gives its singular coordinate
\begin{equation}
 s=\bigl(a_\ell-(\ell+1)F_\ell\bigr)Q
 \label{rd:singular-formula}
\end{equation}
up to the uniformly bounded descent torsion already described.
The Frobenius polynomial yields the exact operator identity
\begin{equation}
 (F_\ell-a_\ell)
     \bigl(a_\ell F_\ell-(\ell+1)\bigr)
       =a_\ell-(\ell+1)F_\ell.
 \label{rd:exact-frobenius}
\end{equation}
Thus \(s=(F_\ell-a_\ell)f\); when \(a_\ell\) vanishes at the
retained torsion precision, this is \(s=F_\ell f\).  Taking the
Frobenius limit gives \(s=C_\tau f\).

If $\gamma$ is the chosen rational Frobenius lift, its image in the
generalized-dihedral ring-class Galois group lies outside the abelian
subgroup, so $\gamma^2$ acts trivially on the full upper ring-class field.
The same lift therefore computes the descended finite Kummer coordinate
from the reduced point over the common residue field $\mathbf F_{\ell^2}$.
For the new finite coordinate, reduction of the derivative point is
\begin{equation}
 \sum_{i=1}^{\ell}iF_\ell X
       =\frac{\ell(\ell+1)}2F_\ell X.
 \label{rd:half-factor}
\end{equation}
The finite Kummer coordinate consequently vanishes modulo \(2^a\)
if \(v_2(\ell+1)\ge a+1\).  The additional \(1\) is essential at
\(2\).  It is imposed in the choice of the next precision, and does
not multiply the class by a new denominator.  Bounded torsion
ambiguities disappear on passing from \(E[2^a]\) to \(E[2^m]\) with
\(a-m\ge e\).  Choose \(a\) and the Frobenius precision successively
faster than \(m\).

The calculation is compatible with every translate in the transfer:
transport the place and its Frobenius before applying that translate.
The scalar weight on the lower and upper class is identical.  Applying
the same calculation at each \(\ell\in I\) proves the earlier
singular conditions as well.

The local formulas now give the switch. To use it in the rank-reduction
argument, we must retain more than its cohomology class: the transferred
target cycles, local lifts and degree-two cochain identities must extend
the original finite diagrams. We record this compatibility next.

Here the transfer itself is integral at the cochain level.
For a finite Artin coefficient ring \(A\), let
\(\Delta=\operatorname{Gal}(H_{cr}/K)\), through which the
finite scalar and quadratic characters factor.  Shapiro's
cochain map and the coefficient homomorphism
\[
 A[\Delta]\longrightarrow A,\qquad
 [\sigma]\longmapsto\chi_e(\sigma)\Psi(\sigma)
 \]
send the descended class to the desired twisted class.
For a universal quadratic family, replace the quadratic
value by its group element in \(A[G]\), also including the
fixed base character.  Choose inverse actions on both
sides when using that Shapiro convention.  Explicitly,
these are the finite coset sums with the indicated scalar
weights; there is no division by \(\#\Delta\).
These maps and their degree-zero and degree-two compatibility formulas
are used at this finite stage, before passage to a bounded target model.
Naturality of restriction and Kummer corestriction
supplies the local lifts. At \(2N\) the scalar character
is trivial, so these are precisely the Kummer conditions
of the indicated quadratic twist.

More explicitly, in the universal case let
$\nu:\Delta\to G$ be the additional quadratic character.
Perform the coefficient map
\[
 A_i[\Delta]\longrightarrow A_i[G],\qquad
 [\sigma]\longmapsto
 \chi_{h_0}(\sigma)\Psi_i(\sigma)[\nu(\sigma)]
\]
and the finite coset sum before any character evaluation. This gives
one cycle over $A_i[G]$, with its local lifts. Every sign evaluation
$A_i[G]\to A_i$ sends this cycle and these lifts to the corresponding
weighted construction, because it commutes with finite sums and
differentials. After contraction, compare an evaluated target model
with its already chosen scalar model using their maps to the same
original target cochains. Retain these comparison matrices and their
contraction homotopies for all the finitely many characters of this
fixed $G$. No character idempotent or division by a class-group
order is used. At $t=0$ the split-prime norm multiplier on this universal
cycle is
\[
 a_{r_i}-2\chi_{h_0}(r_i)[\nu(\Fr_{\mathfrak r_i})].
\]
The two conjugate Frobenius values of $\nu$ agree because each
quadratic character comes from $\Q$. Evaluation at the character
corresponding to $h$ therefore gives exactly
$a_{r_i}-2\chi_h(r_i)$, with the central boundary witness described
above.

Finally contract only the resulting target cycles, local lifts and
boundary witnesses into the models over $K$. Retain the target
comparison maps, the split local Fox complexes, their zero isotropy
homotopies, and their duality and switch identities in degrees zero,
one and two. The proof of Lemma~\ref{rl:finite-limits} shows explicitly
how to transport these identities and their witnesses. In particular
the unbounded cochains over $H_{cr}$ and the full Shapiro maps are not
objects or maps of this limiting diagram.
When a further derivative prime is selected, extend the
same original finite-stage diagram and its contraction
maps, using the same inner ultrafilter.  Do not construct
an unrelated cohomologically equivalent old complex.
The earlier evaluations, including all degree-two
cocycle identities, are consequently retained by the
same contraction maps after each new switch. For a fixed finite set of
switches and fixed $G$, the target rank bounds above apply.
Lemma~\ref{rl:finite-limits} therefore applies to these extended
target diagrams at their common retained precision.
Their limiting classes give
\eqref{eq:h:finite-singular}.  The displayed local formulas agree with the
calculation underlying \cite[Proposition~1.7.4]{Howard2012}; Howard's
odd-prime global hypotheses are not being asserted at \(2\).
\end{proof}

\begin{lemma}[Preservation of the determinant coordinate]
\label{lem:h:determinant-switch}
Let $\mathcal O$ be a height-one discrete valuation ring at which the
complexes and switching conditions are self-dual. Suppose the generic
Selmer group for a set $I$ has rank one and is spanned by $P_I$, and
$v_0=f(P_I)$ is nonzero at the next derivative prime. Then the switched
generic group has rank one and is spanned by $P_{I\cup\{l\}}$, and
\begin{equation}\label{eq:h:switch-equality}
 v_{\mathcal O}(u_I)=v_{\mathcal O}(u_{I\cup\{l\}}).
\end{equation}
This assertion does not require a primitive localization in either
integral plane.
\end{lemma}

\begin{proof}
Compare the strict condition and the relaxed condition of
Lemma~\ref{lem:h:switching-pairing}. Their complexes are dual. Over the
fraction field, the image of relaxed $H^1$ in $F\oplus S$ is
self-annihilating by Poitou--Tate, hence two-dimensional. Strict $H^1$ is
zero because the original rank-one group localizes nontrivially. The
relaxed image meets $F$ in the line spanned by $v_0$ and, by
Equation~\eqref{eq:h:finite-singular}, meets $S$ in the nonzero line
spanned by $C_\tau v_0$. Its intersection with either plane is therefore
exactly a line. This proves the rank and spanning assertions.

For the integral determinant comparison let $A=P_I$ and
$B=P_{I\cup\{l\}}$ be a field basis of relaxed $H^1$. The dual basis
$A^\vee,B^\vee$ belongs to strict $H^2$. In the triangle adding $F[-1]$
to the strict complex, the map $F\to H^2_{\mathrm{str}}$ pairs a vector
$x$ against the relaxed localization. By
Equations~\eqref{eq:h:cross-pairing} and
\eqref{eq:h:finite-singular}, it is a unit times
$\langle x,v_0\rangle B^\vee$. The remaining $H^2$ quotient maps to
$A^\vee$. Choose $x$ mapping to $B^\vee$. Although $v_0$ and $x$ need
not be integral separately, the volume of the ordered basis $(v_0,x)$
of $F$ is a unit of $\mathcal O$, precisely because its symplectic
pairing is a unit. For $S$ use $C_\tau v_0$ and a preimage of
$A^\vee$; its volume is again a unit. The two determinant tensors thus
map, up to units, to the same tensor in the strict determinant, using
$A^\vee\wedge B^\vee$. All triangles and plane lattices in this
comparison are integral. Their determinant isomorphisms consequently
give Equation~\eqref{eq:h:switch-equality}, including any torsion or
localization-index contribution.
\end{proof}

\subsection{Rank reduction and nonnegative divisors}\label{sec:h:rank-reduction}

The derivative-prime reduction follows the Kolyvagin-system argument of
\citet[Lemmas~1.5.3 and~1.6.4]{Howard2012}. We carry its switches
through integral Selmer complexes and preserve their determinant
coordinate. The local calculation and descent above supply the
characteristic-two pairing and one fixed multiplier; the argument
below supplies the uniformity as the conductor varies.

We now prove the first two bounds in Equation~\eqref{eq:h:divisor-bounds}.
A height-one specialization can have larger Selmer dimension than the
generic rank-one line. Local switches will remove that excess dimension
while preserving the determinant valuation. At test dimension one,
integrality is immediate. At the coefficient prime $(2)$ this argument
uses residual irreducibility and therefore applies here to genus data.

Use the sequence groups $\Gamma_K\subset\Gamma$ of
Lemma~\ref{rl:evaluation}, and write $\mathcal H$ for the exact
kernel of Lemma~\ref{rl:kernel}. Over the relevant test field, put
$W=\operatorname{Ind}_{\Gamma_K}^{\Gamma}(T\chi_e\Psi)$.
Lemma~\ref{rl:chebotarev-evaluation} realizes finite localization as
the projection of $(1+\tau)b(u')$, for $u'\in\mathcal H$, to one
induced summand. Restriction to $\mathcal H$ is injective by
Lemma~\ref{rl:kernel}, and evaluation there is additive and
Galois-equivariant. The image of $1+\tau$ is a copy of the
two-dimensional summand, also in characteristic two.

\begin{lemma}[Rank reduction at a height-one prime]
\label{lem:h:height-one-reduction}
Suppose a primitive instance has generic cohomology of rank at most one
in degrees one and two, zero elsewhere, and $P$ spans $H^1$ when it is
nonzero. Then the first line of Equation~\eqref{eq:h:divisor-bounds}
holds. These hypotheses hold in every simple-zero instance and also at
the base of Theorem~\ref{thm:heegner-interpolation}. The second line holds in
genus data with irreducible $E[2]$, with a bound independent of the
number of active or derivative primes.
\end{lemma}

\begin{proof}
There is nothing to prove if $u=0$. Otherwise the generic rank is one
and $P$ is nonzero. Localize, and if convenient complete, at a
height-one prime. Test its
residue field; at $(2)$ do this argument only in genus data with
irreducible residual representation. Suppose first that $\Psi^2\ne1$
on the field. Then $W$ is irreducible. Inertia at $r$ distinguishes its
two induced summands. On one summand, irreducibility of $T$ over $K$
comes from open image or the CM normalizer in characteristic zero, and
from the assumed residual irreducibility at $(2)$. In the latter case
we have omitted the fields meeting $\Q(E[2])$ nontrivially. The cyclic
cubic residual case remains irreducible over $\F_2((t))$, with a
possibly quadratic commuting field. Constant Tate matrix images are
realized on the sequence group; scalar character twists do not change
this irreducibility.

Let the test Selmer $H^1$ have dimension $s\ge2$. Include the reduction of
a primitive generator of the generic free $H^1$ lattice among two
independent classes. We can arrange that their localization has rank
two. To see this, if the classes are independent over the commuting
field of $W$, simplicity and the injective restriction in
Lemma~\ref{rl:kernel} imply that the field span of their joint
evaluations on $\mathcal H$ is the whole of $W\oplus W$. Otherwise
a nontrivial equivariant linear relation would vanish on restriction,
and hence on the original classes. Projection by $1+\tau$ therefore
spans all pairs in the two-dimensional localization space. Some actual
evaluation pair has nonzero determinant: a homogeneous quadratic
vanishing on an additive group vanishes on its field span, since its
polar form vanishes on every pair from the group. This argument is
valid in characteristic two as well. If the classes are dependent over
the commuting field but independent over the ground field, projection
commutes with that field action on the graph of $\tau$. At the residual
prime~$(2)$, in the cyclic cubic case, this commuting field is
$\F_4((t))$: rational complex
conjugation has trivial residual image, so the commuting-field action
is the same on the two induced summands. If the second class is $\alpha$ times the first, with
$\alpha$ outside $\F_2((t))$, a nonzero projected evaluation $v$ gives
the ground-field-independent pair $v,\alpha v$. Thus this case also
gives rank-two localization.

Choose the derivative prime realizing that evaluation. Its localization
is nonzero on the primitive generic generator. The relaxed image on the
test fiber is self-annihilating of dimension two, and contains all of
$F$ because the finite localization has rank two. It is therefore $F$.
Switching to $S$ lowers test $H^1$ dimension by two. Generically,
Lemmas~\ref{lem:h:derivative-classes} and
\ref{lem:h:determinant-switch} keep rank one, a nonzero class, and the
same determinant valuation.

It remains to treat $\Psi^2=1$, which occurs here only in
characteristic zero. Extend $\Psi$ across $\tau$ and split into the
two conjugation signs. Conjugation on the cochains includes
$1+t\mapsto(1+t)^{-1}$, which is trivial on this fiber. The fixed
Kummer conditions and relaxed conditions are preserved. At previously
switched places the split local model has diagonal conjugation action
on $f,s$, so its plane conditions are preserved too. The induced
representation now splits into two nonisomorphic absolutely irreducible
twists. Finite and singular localization each have one line per sign.
Within a sign the pairing is a symmetric hyperbolic pairing of those
two isotropic lines; the signs are orthogonal.

If both Selmer signs occur, the restriction argument just used gives
rank-two localization by choosing one class of each sign, with a
nonzero component of the primitive generator among them. If only one
sign occurs and $s>1$, choose rank-one localization nonzero on that
generator. The relaxed image has a self-annihilating line in each
sign. In the sign hit this is the finite line; in the other it must
be the singular line, since a symmetric hyperbolic plane in
characteristic zero has exactly these two isotropic lines. The switch
preserves the total dimension but introduces the other sign: the old
strict kernel has dimension $s-1\ge1$ and the new singular line
contributes the second sign. A subsequent rank-two switch reduces
dimension. Iteration reaches test rank one; it cannot pass below one,
because the generic rank-one class has been retained throughout.

At test rank one the minimal self-dual complex over the DVR has one
free term in each of degrees one and two and zero differential; there
are no invariants. An integral class and its integral dual functional
therefore give an integral determinant coordinate. Reversing the
switches gives the first assertion. At $(2)$ the initial fixed descent
and local multiples cost a bounded valuation, independent of every
subsequent switch. This proves the second assertion with the stated
uniformity.
\end{proof}

\subsection{Uniformity without residual irreducibility}
\label{sec:h:uniform-ring-bound}

Only the ring-data bound at $(2)$ remains. Residual irreducibility is
not assumed, so the preceding argument on that residue field does not
cover every instance. We first bound the model sizes and retain evaluation through a second
limit. These constructions turn a hypothetical
sequence of unbounded negative determinant valuations into a
characteristic-zero outer limit. There the retained evaluation maps
permit rank reduction and force a bounded denominator, giving a
contradiction.

\begin{lemma}[Uniform residual sizes]\label{rl:residual-bound}
Suppose $\omega(\operatorname{disc}K)\leq A$, the local twisting classes
at $2N$ belong to a fixed finite list, and all active conductor primes
are split with odd Frobenius exponent of the deformation character.  After localization and completion at
$(2)$, the minimal free ranks of the Selmer complexes are bounded in
terms of $E,A$, and the number of retained derivative primes.  They do
not depend on the number of active conductor primes.
\end{lemma}

\begin{proof}
Reduce first over $\mathbf F_2[[t]]$.  The quadratic character is now
trivial.  Removing the active conductor primes from the permitted
ramification set changes the global complex by their singular local
complexes.  These, and their unramified local complexes, are acyclic
after inverting $t$, by Lemma~\ref{rl:local-acyclic}.  The resulting
smaller-support complex, including its comparison maps, can be
constructed at the original Artin stages before taking the limit.
Its support consists of $2N,r$, and the retained derivative primes.
Its $\mathbf F_2((t))$ cohomology dimensions are bounded by those at
$t=0$: choose a minimal finite free model over $\mathbf F_2[[t]]$.

Here is a uniform bound for the latter residual cohomology.  Put
$F=\mathbf Q(E[2])$, a fixed field.  Then $FK/F$ is at most quadratic.
Enlarge the allowed set of primes of $F$ by those ramified in $K$.
Its cardinality is bounded by a constant depending on $E,A$ and the
number of derivative primes.  Kummer theory bounds the number of
generators of the maximal pro-$2$ Galois group over $F$ unramified
outside this set, using the fixed unit group, the fixed $2$-class group,
and the allowed valuations.  The subgroup defining $FK$ has index at
most two; the elementary Schreier bound bounds its generator number.
Every finite Galois $2$-extension of $FK$ with this ramification has
Galois closure of $2$-power degree over $F$, so this subgroup indeed
controls the required pro-$2$ extensions of $FK$.

Over $FK$ the residual Tate module is trivial.  Restriction and
inflation--restriction therefore bound the residual $H^1$ over $K$.
The $H^0$ bound is immediate, and the global Euler characteristic or
global duality bounds $H^2$.  The local terms at the remaining bounded
set have bounded dimensions, with degrees of the local fields at most
two over the fixed rational local fields.  At $2N$ there are only
finitely many local quadratic twisting classes.  The same bounds
apply to the Selmer cone.  Lemma~\ref{rl:finite-model} now bounds the
minimal ranks. Determinant valuations will require a separate argument.
\end{proof}

\subsubsection{Evaluation through the outer limit}

The preceding evaluation argument concerns one inner construction over a
fixed $K$. The ring bound must also be uniform when $K$ and the conductor
vary. When all active primes split and $\omega(k)$ is bounded,
Lemma~\ref{rl:residual-bound} bounds the model sizes after any fixed
number of derivative operations. The next proposition shows that passing to a second
limit of such models preserves detection, including classes represented
only by approximate cycles before that limit.

\begin{proposition}[Integral evaluation through the outer limit]
\label{rl:outer-evaluation}
Let $C_j$ be any family of the preceding limiting complexes, completed
over $\mathcal O_j=\widehat{\Lambda_{(2)}}$, possibly for different
quadratic fields $K_j$.  Suppose their minimal ranks are bounded.
Assume a common bounded cohomological degree interval. Keep the
actual evaluation matrices and transport them to the minimal models
through integral contractions.  For a nonprincipal
ultrafilter $\mathcal U$, put
\[
 \mathcal O_\infty=
 \left(\prod_{\mathcal U}\mathcal O_j\right)/\mathfrak a,
 \qquad
 \mathfrak a=\{(x_j):
    \{j:v_2(x_j)\geq n\}\in\mathcal U\text{ for every }n\}.
\]
Then $\mathcal O_\infty$ is a DVR with uniformizer $2$ and a
characteristic-zero fraction field $L_\infty$. The outer $H^1$ over
$L_\infty$ injects into crossed classes on the inner quadratic sequence
group product, or its extension by one simultaneous complex conjugation
for the induced representation.  Its restriction to the product
of their exact common kernels is injective.  All arbitrary-sequence
degree-two evaluation identities survive.  These conclusions also
hold after retaining any fixed finite collection of additional
integral diagrams.
\end{proposition}

\begin{proof}
A nonzero element of the quotient has a representative whose
valuations are bounded on an ultrafilter-large set.  A finite
partition then makes its valuation equal to a single integer $n$
on such a set.  It is $2^n$ times a unit.  Products add these finite
valuations, so the quotient is a domain, and every nonzero ideal has
an element of smallest valuation.  It is consequently a DVR.  The
integer $2^n$ is nonzero for every $n$, so its fraction field has
characteristic zero.

For each $j$, all-sequence evaluation maps originally have entries
in $\Lambda$.  Extension to $\mathcal O_j$ and composition with an
integral contraction still give integral matrices.  Thus evaluation
on any sequence of inner group elements has an entrywise image in
$\mathcal O_\infty$.  Equations~\eqref{rl:eval-boundary} and
\eqref{rl:eval-cocycle} survive for all such sequences.  In particular,
if an outer cycle is represented by vectors $x_j$ for which
$d_jx_j$ has valuation tending to infinity, the evaluation of
$d_jx_j$ does too: the degree-two evaluation matrix is integral.
Thus an outer cycle gives an exact crossed cocycle even when its
representatives are only approximate inner cycles.

Over the inner residue field $\Omega_j$, choose a detecting list as
in Lemma~\ref{rl:evaluation}; its length is at most
$\dim_{\Omega_j}H^1(C_j\otimes\Omega_j)+\operatorname{rank}M$.
These lengths are bounded.  Include in the integral diagrams the
evaluation map to its two-term complex and a contraction of its cone
to a complex starting in degree one.  Such a contraction uses only
unit pivots over $\mathcal O_j$ and has bounded graded sizes.
All its matrices, including inverses and homotopies, are integral
and pass to the quotient.  Its cone still starts in degree one
after extension to $L_\infty$.  This proves injectivity just as in
Lemma~\ref{rl:evaluation}.

Let $H_j$ be the inner exact kernel. In the product of the quadratic
sequence groups, the central elements of Lemma~\ref{rl:kernel} give
a central element modulo $\prod_jH_j$ with scalar $u^4$. Since
$u^4-1$ remains nonzero in $\mathcal O_\infty$, the same calculation
gives restriction injectivity. For the induced representation first
restrict to the quadratic subgroup, as in Lemma~\ref{rl:kernel}.
Additional finite diagrams pass through the same entrywise operations.
\end{proof}

\begin{lemma}[Irreducibility in the outer limit]
\label{rl:outer-irreducibility}
On the outer characteristic-zero test, $\Psi$ has infinite order and
the four-dimensional representation induced from $T\chi\Psi$ is
absolutely irreducible.
\end{lemma}

\begin{proof}
The inner irreducibility argument requires a uniform finite
witness to survive this limit.  The images of the quadratic subgroups have index at most
two in the fixed compact $2$-adic Tate image.  A compact $2$-adic
analytic group is topologically finitely generated and hence has
only finitely many closed subgroups of index at most two.  Choose
one on an ultrafilter-large set.  Its absolutely irreducible
representation has finitely many matrices spanning
$M_2(\mathbf Q_2)$; the determinant expressing this spanning
property is a fixed nonzero $2$-adic number.  Every one of these
Tate matrices is realized by choices of Galois elements in each
component.  Their accompanying character scalars are units and
do not change the spanning assertion.  Its determinant remains
nonzero in $L_\infty$.  This proves absolute irreducibility in
that field.  Also $1+t$ has infinite order already in every
$\Omega_j$: for a nonzero integer $n$, the first nonzero term of
$(1+t)^n-1$ occurs in finite degree.  It remains a unit modulo
$2$ in every component, so no such power becomes one in the
outer field.

The two induced summands therefore have different inertia scalars,
$\Psi$ and $\Psi^{-1}$, and the induction is absolutely irreducible
as in Lemma~\ref{rl:irreducibility}.
\end{proof}

\begin{lemma}[Prime realization through the outer limit]
\label{rl:outer-chebotarev}
The finite Chebotarev realization of Lemma~\ref{rl:chebotarev-evaluation}
is compatible with Proposition~\ref{rl:outer-evaluation} and with any
fixed finite list of previous derivative operations.
\end{lemma}

\begin{proof}
Apply Lemma~\ref{rl:chebotarev-evaluation} at each finite stage.
For an outer sequence do this separately
for each inner sequence, retaining the relevant finite evaluation
matrices and degree-two identities before taking the outer
quotient.  Approximate inner cycles cause no problem, by the
integrality calculation in Proposition~\ref{rl:outer-evaluation}.
Previously imposed finite diagrams are unchanged.
\end{proof}

\begin{theorem}[Compatibility of the limit constructions]
\label{rl:limit-evaluation-package}
The finite free models, ring-class characters, local conditions and
pairings constructed above can be retained simultaneously. Their
inner limits have the stated height-one acyclicity and all-sequence
cohomology detection. For all-split ring data with bounded
$\omega(\operatorname{disc}K)$, the residual model sizes are bounded
independently of the active conductor support. Their outer limits
retain arbitrary evaluations, exact-kernel injectivity and finite
Chebotarev realization after any fixed number of additional
integral arithmetic operations.
\end{theorem}

\begin{proof}
Use Lemma~\ref{rl:finite-model} at each finite stage and keep the
same contractions and ultrafilter as in Lemma~\ref{rl:finite-limits}.
Perform the finite ring-class transfer and weighting first, retaining
only the resulting bounded target cycles and witnesses as in that
lemma; no limit of the growing Shapiro source is taken.
The local assertions and uniform residual sizes are
Lemmas~\ref{rl:kummer}, \ref{rl:local-acyclic} and
\ref{rl:residual-bound}; the required characters come from
Proposition~\ref{rl:characters}. Lemmas~\ref{rl:evaluation} and
\ref{rl:kernel} give inner detection. Proposition~\ref{rl:outer-evaluation}
and Lemma~\ref{rl:outer-chebotarev} preserve it through the
outer limit and later finite operations. Every retained identity is
an identity of actual integral matrices. The finite derivative classes and local identities of
Lemma~\ref{lem:h:derivative-classes} are among the integral arithmetic
target diagrams retained in these operations, including the central
norm relation before and after sign evaluation.
\end{proof}

\begin{lemma}[Uniform determinant bound in ring data]
\label{lem:h:uniform-ring-determinant}
For simple-zero ring data with every active conductor prime split in
$K$ and $\omega(k)\le a$, the coordinate of
Definition~\ref{def:h:determinant-coordinate} satisfies
\[
 v_{(2)}(u)\ge -C_{E,a},
\]
where $C_{E,a}$ is independent of $h$ and $K$.
There is no residual irreducibility assumption. This is the last
line of Equation~\eqref{eq:h:divisor-bounds}.
\end{lemma}

\begin{proof}
Put $O=\widehat{\Lambda_{(2)}}$. By
Lemma~\ref{rl:residual-bound}, after any fixed number of derivative
switches the minimal $O$-models have bounded size in terms of $E$, $a$
and that number, independently of $h$. The full local complexes at
$r$ and the active split primes are contractible over $O$, and the
switched planes are integrally self-dual. Thus the class and its paired
closed functional have a common bounded clearing exponent, supplied by
the fixed descent multiplier and the retained global duality map.

Suppose the asserted lower bound fails. Choose instances with
$v_O(u)\to-\infty$. Cancel disks over $O$ and take an outer limit of
their bounded free models. Transport the class, closed functional,
evaluations and the retained comparison and switch maps through these
integral contractions, applying the same contractions to their boundary
witnesses. Only the resulting bounded target vectors and matrices enter
the outer diagram; the eliminated list of active-place complexes does
not. Concretely, take the ultraproduct of their
valuation rings and quotient by the ideal of sequences whose valuation
tends to infinity. The resulting ring $\mathcal O_\infty$ is a DVR
with uniformizer $2$: a nonzero sequence has a finite valuation on an
ultrafilter-large set and is that power of $2$ times a unit. Its
fraction field has characteristic zero. All integral matrix identities,
duality diagrams, and the local switch diagrams pass to this limit.

Proposition~\ref{rl:outer-evaluation} gives injective evaluation
on the outer sequence group and on its exact common kernel. This
includes outer cycles represented by approximate inner cycles, since
their coboundaries are evaluated by integral degree-two matrices.
The fixed central scalar remains different from one in the outer
field. Lemma~\ref{rl:outer-irreducibility} gives absolute irreducibility
of the induced representation there: the scalar $1+t$ has infinite
order and distinguishes the two summands.

Each individual minimal $O$-complex is a square complex in degrees
one and two, by self-duality and the absence of residual invariants.
Choose a primitive vector in its generic kernel. Such vectors survive
the outer limit, so outer $H^1$ has rank at least one. If its rank exceeds
one, choose rank-two localization by the absolutely irreducible
evaluation argument in Lemma~\ref{lem:h:height-one-reduction}, making
it nonzero on one of these primitive vectors. The inner Chebotarev construction realizes the corresponding
sequences of derivative primes. At each original Artin stage retain
the old contractions and evaluations, adjoin the new local and
derivative diagrams, and pass through the same inner and outer
ultrafilters. Slower cofinal coefficient precision is permitted;
independent reconstruction of the old diagrams is not. For an ultrafilter-large set of instances the localization is
generically nonzero. Lemma~\ref{lem:h:determinant-switch} preserves
the valuations tending to minus infinity, while the outer rank falls
by two. The new complexes and integral classes obey the same bounds,
with only the now fixed finite number of added derivative primes.

The process must reach outer rank one: ranks are finite and cannot
fall below one while the primitive kernel vector is retained. At this
stage some maximal boundary minor is nonzero over the outer fraction
field. Its valuation is consequently bounded in the individual DVRs
on an ultrafilter-large set. Split that minor over their fraction
fields. The entries of the class and duality maps are integral and
their common initial multipliers are bounded. The formula for their
determinant coordinate, or the elementary-divisor decomposition, then
has a lower valuation bound depending on this minor and the bounded
complex size. The Heegner classes themselves need not survive the outer
limit: a numerator tending to zero only increases its valuation. It is
the primitive kernel vector that preserves rank, and the boundary minor
that bounds the denominators. This contradicts $v_O(u)\to-\infty$, and
proves the uniform estimate.
\end{proof}

\begin{proof}[Proof of Theorems~\ref{thm:heegner-genus-bound} and
\ref{thm:heegner-ring-bound}]
By Lemmas~\ref{lem:h:height-one-reduction} and
\ref{lem:h:uniform-ring-determinant}, the element $u$ has nonnegative
valuation away from $(2)$ and a uniformly bounded negative valuation
at $(2)$. Normality of $\Lambda$ implies that a fixed power of $2$
times $u$ belongs to $\Lambda$. Its central value therefore has the
same lower valuation bound.

In genus data, $s(h_*)$ and the weight of $h_*$ are fixed, $c=1$,
and $k=hh_*$. Equations~\eqref{eq:h:central-determinant} and
\eqref{eq:h:weil-restriction-length} give
Equation~\eqref{eq:h:genus-bound}. The excluded intersection fields
are finite in number, since an irreducible two-division field has
at most one quadratic subfield. In ring data, the sum over $q\mid k$
is bounded by $2\omega(k)$, whereas
$2\sum_{q\mid h}t_q=4w(h)$. The same equations give
Equation~\eqref{eq:h:ring-bound}.
\end{proof}

\subsection{The universal Heegner class}
\label{sec:h:universal-class}

For interpolation we combine the non-$(2)$ divisor estimate and the
central valuation formula with a universal binary family. The
characterwise coordinates must have a uniformly controlled clearing
factor. We first remove the
norm and local determinant factors belonging to primes unused at a
character, then eliminate the new-prime blocks before choosing that
clearing factor.

Recall the binary coefficient rings:
\[
 G\simeq(\Z/2\Z)^b,\qquad R=\Lambda[G],\qquad
 O=\widehat{\Lambda_{(2)}}.
\]
We work with the universal quadratic character over $R$ for the
family in Theorem~\ref{thm:heegner-interpolation}. At a new prime $q$
let $g_q\in G$ be its inertia bit: at a character $x$, it evaluates
to $-1$ when $q$ is active and to $1$ when it is unused. Choose the
conductor character of Proposition~\ref{rl:characters} with odd
exponents at all the new split primes. In the universal complex relax
at every new prime and at the old active primes, and take the trace
class $p$ from the conductor containing all new primes, together with
$r$ and the old conductor. As in the proof of
Lemma~\ref{lem:h:derivative-classes}, form its $A_i[G]$-weighted
target cycle at each finite stage, together with its local lifts and
boundary witnesses, before evaluating any sign character. Apply the
limit construction for this fixed $G$ to those target data; the source
ring-class cochains do not enter the limiting diagram. Choose an integral determinant basis and
write $u_{\mathrm{rel}}$ for the characterwise determinant coordinate.

\begin{lemma}[Removing the unused-prime factors]
\label{lem:h:unused-prime-correction}
There is $\eta\in O[G]$ such that
$u=\eta u_{\mathrm{rel}}$ agrees at each character with the primitive
determinant coordinate up to a $\Lambda$-unit. In particular,
$u_x\in\Lambda[1/2]$. If the base product derivative is zero, then
$u_0(0)=0$. At the nonzero vertices,
\begin{equation}\label{eq:h:normalized-vertex-bound}
 v(u_x(0))\le B+O_{E,h_0,a}(1),
\end{equation}
where $a$ bounds $\omega(k)$.
\end{lemma}

\begin{proof}
At an unused prime, the norm relation multiplies the primitive point
class by
\[
 e_q=\begin{cases}
 a_q,&q\text{ inert in }K,\\
 a_q-z-z^{-1},&q\text{ split in }K,
 \end{cases}
\]
where $z=\chi_{h_0}\Psi$ times the universal quadratic scalar at an
oriented Frobenius. The relaxed singular quotient contributes
\[
 l_q=\prod_{v\mid q}
  \det\left(1-q_v^{-1}\eta_v(\Fr_v)\Fr_v\mid T\right),
\]
where $q_v$ is the residue cardinality and $\eta_v$ denotes the full
scalar character. Thus $u_{\mathrm{rel}}$ differs from the primitive
coordinate by $e_q^2/l_q$, up to a $\Lambda$-unit. This follows from
the localization triangle with its integral determinant isomorphism;
the pairing comparison is made over the fraction field. Choose
conjugate-compatible Frobenius lifts in these formulas, taking the
square of rational Frobenius at an inert place, and extend the
expressions universally over $O[G]$.

Both $e_q$ and $l_q$ are units in $O[G]$. At a split prime, the odd
Frobenius exponent gives a nonconstant scalar in $\F_2((t))$; at an
inert new prime, the hypothesis $t_q=0$ gives an odd trace. Moreover,
already in the group ring,
$l_q\equiv e_q^2\pmod2$. In the split case the residual identity is
\[
 (1+a_qz+z^2)(1+a_qz^{-1}+z^{-2})
       =a_q^2+z^2+z^{-2}=(a_q+z+z^{-1})^2.
\]
In the inert case the Frobenius-square polynomial gives $a_q^2$.
Here the group-like quadratic scalars have square one, so the identities
hold before sign evaluation. Consequently
\begin{equation}\label{eq:h:group-ring-correction}
 \eta=\prod_{q\ \mathrm{new}}
  \left(1+\frac{l_q/e_q^2-1}{2}(1+g_q)\right)
       \quad\text{belongs to }O[G].
\end{equation}
Its factor is $l_q/e_q^2$ at unused characters and $1$ at active
characters. It makes the required correction without an idempotent
denominator for each prime.

For all the primitive constituents, generic cohomology has rank at
most one. At nonzero vertices this follows from the simple-zero
assumption. At the base it follows from Selmer corank one; a nonzero
generic class then spans the possible rank-one group. Hence
Lemma~\ref{lem:h:height-one-reduction} applies to every nonzero
primitive determinant coordinate. Its nonnegative valuations away
from $(2)$ imply membership in $\Lambda[1/2]$, by normality. The
same assertion is automatic for a zero coordinate.

If the base derivative vanishes, Equation~\eqref{eq:h:height} makes
the specialized point torsion. Either its generic determinant
coordinate is zero or the specialization is at constant rational
rank one and its degree-one class is zero there. In either case
$u_0(0)=0$. At a nonzero vertex use
Equations~\eqref{eq:h:central-determinant} and
\eqref{eq:h:weil-restriction-length}. Subtracting the two Sha lengths
and $4w(h_x)$ leaves precisely the expression bounded in
Equation~\eqref{eq:h:interpolation-hypothesis}, together with bounded
terms at $k$, the old support and $r$. This gives
Equation~\eqref{eq:h:normalized-vertex-bound}.
\end{proof}

The corrected coordinates have the required upper bounds at nonzero
vertices and vanish at the center if the base derivative vanishes.
To apply binary transfer, we still need $f,d\in O[G]$ with $d_x=f_xu_x$ at every character.
The required controls are a bound on $v(f_0(0))$ and, at each fixed
precision, a bound on the negative Laurent support of every $f_x$. The next lemma obtains those bounds from the
old complex, whose size is bounded, rather than from the full growing
conductor.

\begin{lemma}[A uniform clearing factor for the Heegner family]
\label{lem:h:family-clearing}
For the element $u$ of Lemma~\ref{lem:h:unused-prime-correction},
there are $f,d\in O[G]$ with $d_x=f_xu_x$ at every character,
$f_0\in\Lambda$, and $v(f_0(0))$ uniformly bounded. For each fixed
$M$, the negative Laurent support of every $f_x\bmod2^M$ has a
bound independent of $b$ and of the number of new primes.
\end{lemma}

\begin{proof}
First make the pairing functional integral up to a bounded multiplier.
The class $p$ is integral. Lift $2^{a_0}p$ to the complex strict at
the relaxed places and use integral duality against the relaxed
complex. At $r$ and every new prime the full local complex over
$O[G]$ is contractible: its reduction at the maximal ideal is acyclic,
by the local calculation in Lemma~\ref{lem:h:relaxed-local} and the
root-type alternative. At the old active primes its cohomology is
killed by a bounded power of $2$, because inertia is minus identity.
The old support is fixed. Thus $a_0$ is bounded independently of the
new primes. After inverting $2$, lifting is canonical. We obtain a
closed representative of the paired functional over
$2^{-a_0}O[G]$.

At augmentation, present the complex as the old complex, with no new
primes, extended by the direct sum of their singular complexes.
Each singular complex is a bounded-size block in degrees one and two,
with differential $A_q$ and determinant $l_{q,0}$. The differential is
upper triangular with the old complex as a subcomplex. This graded
free model lifts to $R$: stabilize the finite free models, identify their
augmentations, and lift the corresponding graded changes of basis.
The lifts are invertible because the augmentation ideal lies in the
Jacobson radical. Cancel the additional lifted unit disks before
retaining the old-plus-singular grading. Their augmented cross terms
are zero, so cancellation leaves precisely the displayed augmented
model; the added disks do not enlarge its old part.
The old model has bounded size, since $h_0$ is fixed and
$\omega(k)$ is bounded.
Indeed, the fixed-field Kummer and Schreier argument in
Lemma~\ref{rl:residual-bound}, now retaining the fixed primes of $h_0$,
bounds the residual global and local dimensions uniformly for
$\omega(k)$ bounded. This size bound is separate from the central
torsion-length bound used below.

Over $O[G]$ the entire new-prime block is invertible. Schur elimination
leaves a complex $D$ of the bounded old size, with augmentation exactly
the old complex. For each $A_q$, the order of its residual determinant
at $t=0$ is bounded uniformly. At a split prime this follows from the
odd Frobenius exponent and the degree-two Frobenius polynomial; at a
root-type prime the constant term is already nonzero. Every sign
evaluation of the lifted block is its augmentation plus a matrix
divisible by $2$ with power-series entries. Therefore
Lemma~\ref{lem:h:bounded-poles} bounds the negative Laurent support of
the entries of every $D_x$ modulo each fixed $2^M$, independently of
the number of blocks.

The integral central specialization of the old complex is controlled
by the ordinary base Kummer complex, the fixed old conductor primes,
and the bounded singular contributions at $r$. Its rational $H^1$ and
$H^2$ both have dimension one. Its torsion lengths are bounded by the
finite Selmer quotient assumed in
Theorem~\ref{thm:heegner-interpolation}, together with bounded global
torsion and these local contributions. Choose its elementary-divisor
bases at the center and lift. Apply case \textup{(i)} of
Lemma~\ref{lem:h:pivot-clearing} to $D$, the projected cycle and
paired functional. The eliminated large block contributes an
$O[G]$-unit to the determinant, and the multiplier in
Equation~\eqref{eq:h:group-ring-correction} belongs to $O[G]$.
Include both in the allowed factor $\eta$ of that lemma, so neither
enters the clearing factor $f$. The lemma gives $d=fu$, a bounded
central valuation for $f_0$, and the asserted pole bounds for $f_x$.
\end{proof}

\begin{proof}[Proof of Theorem~\ref{thm:heegner-interpolation}]
The primitive ring-class product has odd functional-equation sign
under the standing Heegner conditions. Suppose the base product does
not have a simple zero. Its first derivative then vanishes. By
Lemma~\ref{lem:h:unused-prime-correction}, $u_0(0)=0$, whereas every
nonzero vertex has central determinant valuation bounded by one
constant $B_1$ independent of the family. By
Lemma~\ref{lem:h:family-clearing}, the hypotheses of
Lemma~\ref{lem:h:boolean-transfer} hold with uniform central and pole
bounds. Choose its fixed precision larger than $B_1$ plus twice the
central bound, and then take the cube dimension sufficiently large.
It supplies a nonzero vertex whose central valuation exceeds $B_1$,
a contradiction. Thus the base product has a simple zero.
\end{proof}
